\documentclass[ECM\space ReStud.tex]{subfiles}

\begin{document}

\FloatBarrier

\startonlineappendix

% appendix numbering reset counting of tables
%\counterwithin{figure}{section}
%\counterwithin{table}{section}
%\renewcommand{\thefigure}{\Alph{section}\arabic{figure}}
%\renewcommand{\thefigure}{\Alph{section}\arabic{table}}


%\setcounter{table}{0}
%\renewcommand{\thetable}{A\arabic{table}}

%\setcounter{figure}{0}
%\renewcommand{\thefigure}{A\arabic{figure}}


%\addcontentsline{toc}{section}{Appendices}
%\renewcommand{\thesubsection}{\Alph{subsection}}

%\section{Online Appendix}

\subsection{Appendix Figures}

\begin{figure}[htbp]
    \centering
    \includegraphics[width=5in]{figures/reported-social-perception-1.pdf}
\caption{Social Approval and Disapproval at Baseline}
\label{fig:praise-stigma-plot}
\floatfoot{\textit{Notes:} The figure reports baseline responses to vignette questions about four behaviors. For each vignette, respondents stated whether they would praise someone who performs the behavior (left panel) and whether they would look down on someone who does not (right panel). Bars report the share answering ``Yes'' and ``No.'' All 1{,}995 baseline adults answered the adult deworming and child immunization vignettes. The latrine and church attire vignettes were randomly assigned to half the sample each.}
\end{figure} %There are no significant differences across experimental arms (see Table~\ref{tab:baseline-social-norms}).

\newpage

\begin{figure}[htbp]
    \centering
    \includegraphics[width=4.5in]{figures/dist-plot-1.pdf}
    \caption{Distance to Assigned Point of Treatment}
    \label{fig:distance}
    \floatfoot{\textit{Notes:} Each panel plots kernel densities of realized community-centroid-to-PoT distances (in meters), separately for communities classified as Close (at most 1.25 km) and Far (more than 1.25 km); the vertical dashed line marks the 1.25 km cutoff. Panels correspond to the four incentive arms (Control, Ink, Calendar, Bracelet). Across all 144 study communities, mean realized distance is 1.02 km greater in Far than in Close communities. Close communities have realized distances of at most 1.25 km, while Far communities have realized distances above 1.25 km. Kernel density smoothing can make the plotted Close and Far distributions overlap around this cutoff and extend beyond the observed distance ranges.}
\end{figure}


\begin{figure}[htbp]
    \centering
    \includegraphics[width=4.5in]{misc-figures/all-clusters-map-1.pdf}
    \caption{PoT Selection and Catchment Areas}
    \label{fig:sites}
    \floatfoot{\textit{Notes:} Black crosses (+) denote selected points of treatment (PoTs), while the grey regions indicate the non-overlapping catchment circles from which we selected communities to target. Circle colors indicate counties (Busia red, Siaya blue, Kakamega green). Clusters and PoTs were selected using a geographically constrained procedure that limits overlap across neighboring clusters. Appendix~\ref{sec:site-selection} provides full details.}
\end{figure}

\newpage


\begin{figure}[htbp]
    \centering
    \includegraphics[width=3.3in]{misc-figures/signal_bracelet.png}
    \caption{Bracelet}
    \label{fig:bracelet}
    \floatfoot{\textit{Notes:} Green silicone bracelet distributed by CHVs at the point of treatment. The Swahili text reads ``Tibu minyoo: boresha afya ya jamii yako'', translated as ``Treat worms: improve the health of your community''.}
\end{figure}

\begin{figure}[htbp]
    \centering
    \includegraphics[width=3.3in]{misc-figures/calendar.png}
    \caption{Calendar}
    \label{fig:calendar}
    \floatfoot{\textit{Notes:} Wall calendar distributed at the point of treatment and used as an active comparison item for public observability. The calendar contains no deworming message.}
\end{figure}


\clearpage
\newpage

\begin{figure}
    \centering
    \includegraphics[width=5in]{figures/appendix/sensitization.png}
    \caption{Pre-Campaign Community Outreach Script}
    \label{fig:sensitization}
    \floatfoot{\textit{Notes:} This figure shows the standardized script that CHVs followed during household visits before the deworming campaign. The script explained that adults may have worms even without symptoms, described the individual and community benefits of regular deworming, and encouraged residents to remind family members and neighbors about the campaign. CHVs completed the blank fields with the locally relevant campaign date and PoT location and, where applicable, the item or signal assigned to the community.}
\end{figure}

\clearpage
\newpage


\begin{figure}[htbp]
    \centering
    % First Row with two figures
    \begin{subfigure}[b]{0.5\textwidth}
        \centering
        \includegraphics[width=2.2in]{misc-figures/flyer_bracelet.png}
        \caption{Bracelet Flyer}
        \label{fig:sub_flyer_bracelet}
    \end{subfigure}%
    ~ %add desired spacing between images, e. g. ~, \quad, \qquad, \hfill etc.
      %(or a blank line to force the subfigure onto a new line)
    \begin{subfigure}[b]{0.5\textwidth}
        \centering
        \includegraphics[width=2.2in]{misc-figures/flyer_calendar.png}
        \caption{Calendar Flyer}
        \label{fig:sub_flyer_calendar}
    \end{subfigure}

    % Second Row with two figures
    \begin{subfigure}[b]{0.5\textwidth}
        \centering
        \includegraphics[width=2.2in]{misc-figures/flyer_ink.png}
        \caption{Ink Flyer}
        \label{fig:sub_flyer_ink}
    \end{subfigure}%
    ~ %add desired spacing
    \begin{subfigure}[b]{0.5\textwidth}
        \centering
        \includegraphics[width=2.2in]{misc-figures/flyer_control.png} % Adjusted width for alignment
        \caption{Control Flyer}
        \label{fig:sub_flyer_control}
    \end{subfigure}
    \caption{Flyers}
    \label{fig:flyers}
    \floatfoot{\it{Notes: }\rm This figure displays the four one‑page flyers that CHVs handed out during the pre‑deworming information visits. Each flyer—printed in Swahili—announces the free adult deworming campaign, leaves blank fields for the locally customized site (“\textit{Wapi}”) and date (“\textit{Lini}”), and, except in the Control arm, depicts the incentive offered upon deworming (bracelet, calendar, or ink mark).  Apart from the incentive panel, wording and layout are identical across flyers to preserve comparability.}

\end{figure}



\clearpage
\newpage

\begin{figure}[htbp!]
    \centering
    \includegraphics[width=0.9\linewidth]{figures/gpt-reason-plot.pdf}
    \caption{Reasons for the Observability or Non-Observability of Respondents' Deworming Status}
    \label{fig:sob-reasons}
\floatfoot{\textit{Notes:}
The figure summarizes respondents' explanations for why named peers would or would not know the respondent's own deworming status. When a respondent stated that a named peer knew whether the respondent had dewormed, we asked: ``Why do you think they would think that [you came or did not come for deworming]?'' Verbatim answers were classified with OpenAI's GPT-4 into six substantive categories: (i) direct observation (``He saw me at the PoT,'' ``We went together,'' ``We never met during the campaign''); (ii) communication (``I told them,'' ``We talked about it at the shop,'' ``I told her I will not go''); (iii) social proximity (``She is my neighbor,'' ``We are close friends,'' ``We stay far from each other''); (iv) item/signal cues (for example, seeing the bracelet, ink mark, or calendar); (v) personal characteristics (``They know I take health seriously,'' ``I love health things,'' ``I was pregnant/sick''); and (vi) other. Each stacked bar corresponds to one item/signal arm (Control, Ink, Calendar, Bracelet) in either Close (left four bars) or Far (right four bars) communities. Colored segments show the share of coded reasons in each category. The gray segment labeled ``Don't Know'' shows the fraction of named peers whom the respondent believed did not know the respondent's deworming status. The analysis includes 1{,}512 respondents with classifiable reason responses.}
\end{figure}

\clearpage
\newpage


\begin{figure}[htbp]
    \centering
    \includegraphics[width=4.8in]{misc-figures/takeup-np-rf-fit.pdf}
    \caption{Take-up across Distance by Item/Signal Arm}
    \label{fig:np-rf}
    \floatfoot{\textit{Notes:} This figure plots nonparametric fits of deworming take-up on distance to the point of treatment, estimated separately by assigned item/signal arm: Control, Calendar, Ink, and Bracelet. Each dot is a community level mean among monitored adults. Solid curves are penalized spline estimates of the conditional expectation function. Distance is measured from the community centroid to the assigned point of treatment in kilometers.}
\end{figure}


\clearpage
\newpage

\begin{figure}[htbp]
    \centering
    \includegraphics[width=5in]{figures/sms-TE-by-dist-incentive.pdf}
    \caption{Item/Signal Treatment Effects by SMS Condition}
    \label{fig:te-by-sms}
    \floatfoot{\textit{Notes:} The figure plots item/signal arm treatment effects relative to no-item Control separately for enrolled Social Information SMS recipients and SMS control phone owners. Reminder Only recipients are excluded because the Reminder Only treatment was implemented only in no-item Control communities. Estimates are shown for the full sample and separately for Close and Far communities. Social Information recipients received both a reminder and information about current local take-up, so comparing treatment effects across SMS status provides an auxiliary test of reminder and social learning channels. The analysis includes 2{,}635 Social Information recipients and 3{,}813 SMS control phone owners.}
\end{figure}


\clearpage
\newpage


\begin{figure}[htbp]
    \centering
    \includegraphics[width=4.2in]{misc-figures/dydc-derivative-positive-check.pdf}
\caption{Posterior Distribution of the Cutoff Response to Distance}

    \label{fig:dydc-robust-check}
\floatfoot{\textit{Notes:}
For every posterior draw of the structural parameters and for the range of distances used in the paper, the estimated cutoff response to distance $\bigl(\delta-\mu'(d)\Delta(w^\ast(d))\bigr)\big/
\bigl(1+\mu(d)\Delta'(w^\ast(d))\bigr)$
is strictly positive. Thus increasing distance raises the equilibrium cutoff and lowers take-up in all posterior draws considered.}
\end{figure}


\clearpage
\newpage


\begin{landscape}
    \begin{figure}
            \centering
            \includegraphics[width=10in]{optim-figures/fit105/panel-scenarios-compare-optimal-allocation-plot-distconstraint-3500.pdf}
            \vspace{1pt} % Adjust the vertical spacing between figure and caption
            \caption{Optimal Placement of Treatment Sites (with Experimental Benchmark)}
            \label{fig:oa-comp-panel}
            \floatfoot{\textit{Notes:}
            Each panel depicts one allocation scenario computed at the component-wise median of 1{,}600 posterior draws of the structural parameters. From left to right: (i) the experimental assignment with every community treated as if it were in the Control arm; (ii) the policymaker's optimal allocation under Control observability; and (iii) the policymaker's optimal allocation under Bracelet observability. Black dots mark the 144 communities; red triangles are funded points of treatment (PoTs); blue triangles are candidate PoTs left unused. Black line segments connect each community to its assigned PoT. In both optimal scenarios, the integer program chooses among 1{,}451 candidate sites subject to a 3.5 km distance cap and preserves the expected number of adults treated under the experimental allocation and Control observability, using census adult-population weights. These representative-parameter allocations use 102 sites under Control and 76 under Bracelet; they differ from posterior summaries of allocations solved separately for each draw.}
    \end{figure}
\end{landscape}

\newpage

\newpage
\FloatBarrier
\subsection{Appendix Tables}\label{sec:appendix-tables}


\begin{table}[ht!]
    \scriptsize
    \caption{Balance for Main Analysis Samples}
    \label{tab:stacked-sample-balance-table}
    \centering
    \begin{threeparttable}
    \resizebox{\textwidth}{!}{%
    \input{tables/stacked-sample-balance-table}
    }
    \floatfoot{\textit{Notes:}
    Panel A reports balance for the no-SMS take-up analysis sample. Panel B reports balance for the no-SMS endline observability sample. The ``Sample Mean'' column reports the mean in the corresponding analysis sample. Columns labeled ``Con'' report the no-item Control mean, with standard errors in parentheses. Columns labeled ``Ink - Con'', ``Cal - Con'', ``Bra - Con'', and ``Bra - Cal'' report differences relative to the indicated comparison group, with \(p\)-values in brackets below. ``Joint Test'' reports the F-test \(p\)-value for equality of means across the four item/signal arms within each distance condition, and in the final column across distance conditions combined. Standard errors are clustered at the community level. County fixed effects are included to reflect stratification. The Distance to PoT row reports realized community-centroid-to-PoT distances in each distance category and item/signal arm.}
    \end{threeparttable}
\end{table}

\begin{table}[ht!]
    \scriptsize
    \caption{Balance by Item/Signal Arm: Far--Close Differences}
    \label{tab:stacked-sample-dist-balance-table}
    \centering
    \begin{threeparttable}
    \resizebox{\textwidth}{!}{%
    \input{tables/stacked-sample-dist-balance-table}
    }
    \floatfoot{\textit{Notes:}
    Panel A reports balance for the no-SMS take-up analysis sample. Panel B reports balance for the no-SMS endline observability sample. The table reports the same variables as Table~\ref{tab:stacked-sample-balance-table}, reorganized to show Far--Close differences within each item/signal arm. For each arm---no-item Control, Ink, Calendar, and Bracelet---the first column reports the Close mean, with standard errors in parentheses. The adjacent column reports the Far--Close difference within that arm, with \(p\)-values in brackets below. Standard errors are clustered at the community level. County fixed effects are included to reflect stratification.}
    \end{threeparttable}
\end{table}


\begin{table}[ht!]
	\scriptsize
	\caption{Monitored Take-up Attrition}
    \label{tab:takeup-attrit}
	\centering
    \begin{threeparttable}
    % \resizebox{\textwidth}{!}{%
    
\centering
\resizebox{\ifdim\width>\linewidth\linewidth\else\width\fi}{!}{
\begin{tabular}{lcc}
   \tabularnewline \midrule \midrule
                                                                         & Treatment     & Treatment $\times$ Distance \\    
   Dependent variable: Missing from monitored take-up sample       & (1)           & (2)\\  
   \midrule
   Ink                                                                   & -0.004        &   \\   
                                                                         & (0.005)       &   \\   
   Calendar                                                              & 0.004         &   \\   
                                                                         & (0.005)       &   \\   
   Bracelet                                                              & 0.006         &   \\   
                                                                         & (0.005)       &   \\   
   Far                                                                   &               & 0.006\\   
                                                                         &               & (0.007)\\   
   Ink $\times$ Close                                                    &               & -0.003\\   
                                                                         &               & (0.007)\\   
   Ink $\times$ Far                                                      &               & -0.006\\   
                                                                         &               & (0.007)\\   
   Calendar $\times$ Close                                               &               & 0.003\\   
                                                                         &               & (0.008)\\   
   Calendar $\times$ Far                                                 &               & 0.004\\   
                                                                         &               & (0.005)\\   
   Bracelet $\times$ Close                                               &               & 0.008\\   
                                                                         &               & (0.006)\\   
   Bracelet $\times$ Far                                                 &               & 0.004\\   
                                                                         &               & (0.007)\\   
   \addlinespace[0.4em]
   Female                                                                & 0.015$^{***}$ & 0.015$^{***}$\\   
                                                                         & (0.004)       & (0.004)\\   
   Age                                                                   & 0.001$^{***}$ & 0.001$^{***}$\\   
                                                                         & (0.0002)      & (0.0002)\\   
   Expected distance to PoT                                              & -0.008        & -0.009\\   
                                                                         & (0.011)       & (0.011)\\   
   Control mean                                                          & 0.091         & 0.091\\  
   County FEs                                                            & Yes           & Yes\\  
   RF controls                                                           & Yes           & Yes\\  
   $H_0$: all displayed treatment coefficients = 0, $p$-value            & 0.063         & 0.181\\  
   $H_0$: Bracelet = Calendar, $p$-value                                 & 0.488         & 0.672\\  
   $H_0$: Bracelet $\times$ Close = Calendar $\times$ Close, $p$-value   &               & 0.373\\  
   $H_0$: Bracelet $\times$ Far = Calendar $\times$ Far, $p$-value       &               & 0.926\\  
   \midrule
   Observations                                                          & 10,794        & 10,794\\  
   R$^2$                                                                 & 0.121         & 0.121\\  
   \midrule \midrule
\end{tabular}}



	% }
    \floatfoot{\textit{Notes:}
    The dependent variable is an indicator equal to one if an individual assigned to administrative take-up monitoring was omitted from the PoT monitoring list because of the SurveyCTO list generation error. The sample consists of the no-SMS take-up analysis sample (\(N=9{,}805\)) plus individuals omitted from monitoring because of this error (\(N=989\)), for a total of 10{,}794 observations.   The Control mean row reports the no-item Control mean. Bracelet, Calendar, and Ink
  rows report differences relative to the no-item Control arm. The first column
  reports pooled treatment effects. The second column reports treatment-by-distance
  effects, with Close and Far treatment cells shown separately. Standard errors, in
  parentheses, are clustered at the community level. County fixed effects and the
  reduced-form controls are included throughout.}
\end{threeparttable}
\end{table}


\begin{table}[ht!]
	\scriptsize
	\caption{Endline Observability Module Missingness}
    \label{tab:endline-obs-attrit}
	\centering
    \begin{threeparttable}
 %    \resizebox{\textwidth}{!}{%
    \centering
\resizebox{\ifdim\width>\linewidth\linewidth\else\width\fi}{!}{
\begin{tabular}[t]{lcc}
\toprule
 & Treatment & Treatment $\times$ Distance \\
Dependent variable: Missing from observability module & (1) & (2) \\
\midrule
Ink & \makecell[c]{-0.038\\{(0.036)}} & \makecell[c]{} \\
Calendar & \makecell[c]{-0.031\\{(0.037)}} & \makecell[c]{} \\
Bracelet & \makecell[c]{-0.065$^{*}$\\{(0.036)}} & \makecell[c]{} \\
Far & \makecell[c]{} & \makecell[c]{0.007\\{(0.062)}} \\
Ink $\times$ Close & \makecell[c]{} & \makecell[c]{-0.049\\{(0.041)}} \\
Ink $\times$ Far & \makecell[c]{} & \makecell[c]{-0.019\\{(0.061)}} \\
Calendar $\times$ Close & \makecell[c]{} & \makecell[c]{-0.018\\{(0.047)}} \\
Calendar $\times$ Far & \makecell[c]{} & \makecell[c]{-0.045\\{(0.059)}} \\
Bracelet $\times$ Close & \makecell[c]{} & \makecell[c]{-0.091$^{**}$\\{(0.038)}} \\
Bracelet $\times$ Far & \makecell[c]{} & \makecell[c]{-0.036\\{(0.062)}} \\
\addlinespace[0.4em]
Female & \makecell[c]{0.035$^{*}$\\{(0.020)}} & \makecell[c]{0.036$^{*}$\\{(0.020)}} \\
Age & \makecell[c]{0.000\\{(0.001)}} & \makecell[c]{0.000\\{(0.001)}} \\
Expected distance to PoT & \makecell[c]{-0.046\\{(0.037)}} & \makecell[c]{-0.044\\{(0.039)}} \\
\midrule
Control mean & 0.135 & 0.135 \\
County FEs & Yes & Yes \\
RF controls & Yes & Yes \\
$H_0$: all displayed treatment coefficients = 0, $p$-value & 0.270 & 0.158 \\
$H_0$: Bracelet = Calendar, $p$-value & 0.205 & 0.119 \\
$H_0$: Bracelet $\times$ Close = Calendar $\times$ Close, $p$-value &  & 0.039 \\
$H_0$: Bracelet $\times$ Far = Calendar $\times$ Far, $p$-value &  & 0.833 \\
Observations & 1,267 & 1,267 \\
Simulations & 1,000 & 1,000 \\
\bottomrule
\end{tabular}}

	% }
\floatfoot{\textit{Notes:}
The dependent variable is an indicator equal to one if an endline
respondent is missing the observability module. The sample is the
endline observability module (\(N=1{,}141\)) plus the missing, (\(N=126\)). Bracelet, Calendar, and Ink
rows report differences relative to the no-item Control arm, whose
mean is reported in the Control mean row. The first column reports
pooled treatment effects; the second reports treatment-by-distance
effects, with Close and Far treatment cells shown separately.
Standard errors, in parentheses, are clustered at the community
level. County fixed effects and the reduced-form controls are
included throughout. Module assignment was not recorded for
respondents who did not receive the module; we therefore simulate
1{,}000 random assignments of these respondents to the observability
and perceived observability modules and report estimates averaged
across simulations.}

\end{threeparttable}
\end{table}


% \begin{table}[ht!]
% 	\scriptsize
% 	\caption{Endline Observability Module Missingness by Backup and Deworming Status}
%     \label{tab:endline-obs-attrit-backup}
% 	\centering
%     \begin{threeparttable}
%     % \resizebox{\textwidth}{!}{%
%     \input{tables/know-attrition-backup-dworm-table}
% 	% }
%       \floatfoot{\textit{Notes:}
%     The dependent variable is an indicator equal to one if an endline respondent is missing the observability module. The sample is the no-SMS endline sample (\(N=2{,}659\)). The table reports linear probability models with county fixed effects. \(Bracelet \times Close\) and \(Bracelet \times Far\) report the Bracelet arm missingness gap within each distance condition. ``Backup respondent'' indicates that the respondent was surveyed from the backup endline list rather than the primary endline list. Column 2 additionally controls for own deworming status. Standard errors are clustered at the community level.}
% \end{threeparttable}
% \end{table}


\begin{table}[ht!]
    \scriptsize
    \caption{Lee Bounds for Effects on Observability}
    \label{tab:rf-fob-lee}
    \centering
    \begin{threeparttable}
    \resizebox{\textwidth}{!}{%
    
\centering

\centering
\resizebox{\ifdim\width>\linewidth\linewidth\else\width\fi}{!}{
\begin{tabular}[t]{lcccc}
\toprule
\multicolumn{1}{c}{ } & \multicolumn{4}{c}{Lee Bounds} \\
\cmidrule(l{3pt}r{3pt}){2-5}
\multicolumn{1}{c}{ } & \multicolumn{1}{c}{Combined} & \multicolumn{1}{c}{Close} & \multicolumn{1}{c}{Far} & \multicolumn{1}{c}{Far - Close} \\
Dependent variable: Observability & (1) & (2) & (3) & (4)\\
\midrule
Bracelet & \makecell[c]{[0.127, 0.213]\\(0.041), (0.042)} & \makecell[c]{[0.04, 0.143]\\(0.049), (0.061)} & \makecell[c]{[0.239, 0.303]\\(0.067), (0.05)} & \makecell[c]{[0.096, 0.263]\\(0.09), (0.07)}\\
\addlinespace
Calendar & \makecell[c]{[0.025, 0.088]\\(0.048), (0.051)} & \makecell[c]{[0.019, 0.053]\\(0.058), (0.078)} & \makecell[c]{[0.033, 0.132]\\(0.08), (0.065)} & \makecell[c]{[-0.02, 0.113]\\(0.112), (0.087)}\\
\addlinespace
Ink & \makecell[c]{[0.085, 0.17]\\(0.042), (0.048)} & \makecell[c]{[0.03, 0.139]\\(0.055), (0.067)} & \makecell[c]{[0.154, 0.21]\\(0.065), (0.067)} & \makecell[c]{[0.015, 0.18]\\(0.094), (0.086)}\\
\midrule
\addlinespace
Control & \makecell[c]{[0.679, 0.682]\\(0.033), (0.032)} & \makecell[c]{[0.744, 0.747]\\(0.043), (0.043)} & \makecell[c]{[0.597, 0.598]\\(0.049), (0.046)} & \makecell[c]{[-0.151, -0.146]\\(0.065), (0.063)}\\
\addlinespace
Observations & 1,057 & 1,057 & 1,057 & 1,057\\
\bottomrule
\end{tabular}}


    }
    \floatfoot{\textit{Notes:}
    This table replicates the specification in Table~\ref{tab:beliefs-visibility} using Lee bounds to account for differential missingness in the endline observability module across treatment arms. In each cell, square brackets report \([\text{lower bound}, \text{upper bound}]\), with bootstrapped standard errors for each bound in parentheses below. For the Far--Close column, bounds are derived by interval arithmetic:
    \(\text{lb}(\text{Far}-\text{Close})=\text{lb}(\text{Far})-\text{ub}(\text{Close})\) and
    \(\text{ub}(\text{Far}-\text{Close})=\text{ub}(\text{Far})-\text{lb}(\text{Close})\).
    Standard errors for the Far--Close bounds are computed as \(\sqrt{\sigma^2_{\text{Far}}+\sigma^2_{\text{Close}}}\), where \(\sigma_{\text{Far}}\) and \(\sigma_{\text{Close}}\) are the bootstrapped standard errors for the corresponding bounds. The Lee bounds procedure starts from the 1{,}141 observations in the endline observability sample; trimming leaves \(N=1{,}057\) observations.}
    \end{threeparttable}
\end{table}


\begin{table}[ht!]
    \scriptsize
    \caption{Community Outreach Implementation Checks}
    \label{tab:implement-bal}
    \centering
    \begin{threeparttable}
    \resizebox{\textwidth}{!}{%
    \input{tables/implementation-balance-table}
    }
    \floatfoot{\textit{Notes:}
    The table reports endline measures of campaign outreach and information delivery in the no-SMS endline sample. The ``Sample Mean'' column reports the mean in the full sample. Columns labeled ``Con'' report the no-item Control mean, with standard errors in parentheses. Columns labeled ``Ink - Con'', ``Cal - Con'', ``Bra - Con'', and ``Bra - Cal'' report differences relative to the indicated comparison group, with \(p\)-values in brackets below. ``Joint Test'' reports the F-test \(p\)-value for equality of means across the four item/signal arms within each distance condition, and in the final column across distance conditions combined. Standard errors are clustered at the community level. County fixed effects are included to reflect stratification.
    }
    \end{threeparttable}
\end{table}

\begin{table}[ht!]
 \scriptsize
    \caption{Endline Checks on Item/Signal Receipt, Visibility, and Interpretation}
    \label{tab:incentive-check}
    \centering
    \begin{threeparttable}
    \resizebox{\textwidth}{!}{%
    \input{rf-tables/main-specs/incentive-check-tbl}
    }
   \floatfoot{\textit{Notes:}
    The table reports respondent level indicators collected at endline in the no-SMS endline sample, restricted to the item/signal arms Bracelet, Calendar, and Ink. The Ink row reports the Ink-arm mean. The Bracelet and Calendar rows report differences relative to the Ink arm, with standard errors clustered at the community level in parentheses. ``Seen in community'' is defined among respondents in the item/signal arms and equals one if the respondent reports having seen the assigned item or mark in the community. ``Links to deworming'' is defined among respondents who report having seen the item or mark and provide a non-missing meaning response. It equals one if the response mentions deworming, medication, treatment, tablets, drugs, or worms. ``Received assigned item/signal'' is defined among dewormed respondents and equals one if the respondent reports having received the assigned bracelet, calendar, or ink mark at the PoT. ``Still has / visible'' is defined among dewormed respondents and equals one if the respondent still possessed the assigned bracelet or calendar, or displayed a visible ink mark at endline. ``Wearing / displayed'' is defined among dewormed respondents and equals one if the respondent was wearing the bracelet, had the calendar displayed, or displayed a visible ink mark at endline. \(\sym{*}p<0.10\), \(\sym{**}p<0.05\), \(\sym{***}p<0.01\).}
    \end{threeparttable}
\end{table}

\begin{table}[ht!]
  \centering
  \small
  \caption{Social Image Measures by Baseline Externality Knowledge}
  \label{tab:praise-stigma-externality}
  \begin{threeparttable}
    \input{rf-tables/praise-stigma-by-externality-knowledge}
   \floatfoot{
    \textit{Notes:} This table reports how likely respondents are to praise someone who participates in the deworming campaign or to look down on someone who abstains at baseline, by externality knowledge. Each column reports coefficients from community-level OLS regressions where the dependent variable is the share of baseline respondents in the community who (i) say they would praise someone who participates in the deworming campaign (left column) or (ii) say they would look down on someone who abstains (right column). Panel A uses a binary measure of externality knowledge: “Above median knowledge” equals 1 if the community’s baseline externality knowledge exceeds the sample median. Panel B uses the community’s baseline average externality knowledge (``Know worms impose externalities'', see Table~\ref{tab:baseline-bal}). Heteroskedasticity‑robust standard errors are in parentheses. \sym{*} $p<0.10$, \sym{**} $p<0.05$, \sym{***} $p<0.01$.
   }

  \end{threeparttable}
\end{table}


\begin{table}[ht!]
	\small
	\caption{Distance and Travel to Deworming Sites}
    \label{tab:dist-travel-treatment}
	\centering
    \begin{threeparttable}
    %\resizebox{\textwidth}{!}{%
    \input{rf-tables/main-specs/travel-time-tbl}
	% }
    \floatfoot{\textit{Notes:}
    Travel time and walking shares are computed among endline respondents in the no-SMS sample who dewormed and report a non-zero travel time (\(N=1{,}390\), out of \(N=2{,}659\) endline respondents). Distance to PoT is computed from GIS data for the full take-up analysis sample of 9{,}805 adults.}
\end{threeparttable}
\end{table}


\begin{table}[ht!]
	\scriptsize
    \caption{Balance for SMS Assignment and Enrolled Samples}
    \label{tab:sms-balance}
	\centering
    \begin{threeparttable}
    % \resizebox{\textwidth}{!}{%
    \input{tables/sms/sms-balance-table}
	% }
 \floatfoot{\textit{Notes:} Panel A compares phone owners initially sampled for SMS recruitment with phone owners in the primary SMS control group before SMS recruitment. Panel B compares enrolled SMS recipients with the same SMS control phone owners. The table reports balance on available individual level census covariates. Phone ownership is omitted because all individuals in the table own phones. Columns report the SMS control mean, with standard errors in parentheses, and the SMS treatment minus SMS control difference, with \(p\)-values from tests of equality in brackets. Standard errors are clustered at the community level. County fixed effects are included in all regressions. Panel B shows that enrolled SMS recipients differ from SMS control phone owners in age and gender. All SMS treatment-effect specifications therefore control for age and gender.}
\end{threeparttable}
\end{table}


\begin{table}[ht!]
	\scriptsize
    \caption{SMS Enrollment by Distance and Item/Signal Assignment}
    \label{tab:sms-enrollment}
	\centering
    \begin{threeparttable}
    % \resizebox{\textwidth}{!}{%
    \input{tables/sms/sms-enrollment-balance-table}
	% }
\floatfoot{\textit{Notes:} The dependent variable is an indicator equal to one if an individual sampled for SMS recruitment was reached, consented, and enrolled in the SMS service. The sample is restricted to phone-owning adults initially sampled for SMS recruitment. Columns report enrollment rates by distance condition and item/signal arm, and differences relative to the no-item Control arm within each distance condition. Standard errors are clustered at the community level. Joint test \(p\)-values are from F-tests of equality across item/signal arms within each distance condition and across both distance conditions jointly. County fixed effects are included in all regressions.}
\end{threeparttable}
\end{table}

\begin{table}[ht!]
	\scriptsize
	\caption{Effects of Item/Signal Arms and Continuous Distance on the Observability of Deworming}
    \label{tab:rf-fob-tes-cts-dist}
	\centering
    \begin{threeparttable}
    \resizebox{\textwidth}{!}{%
    \centering
\resizebox{\ifdim\width>\linewidth\linewidth\else\width\fi}{!}{
\begin{tabular}[t]{lcccc}
\toprule
\multicolumn{1}{c}{ } & \multicolumn{4}{c}{Reduced Form} \\
\cmidrule(l{3pt}r{3pt}){2-5}
\multicolumn{1}{c}{ } & \multicolumn{1}{c}{Combined} & \multicolumn{1}{c}{Close} & \multicolumn{1}{c}{Far} & \multicolumn{1}{c}{Far - Close} \\
Dependent variable: Observability & (1) & (2) & (3) & (4)\\
\midrule
Bracelet & \makecell[c]{0.158***\\{(0.037)}} & \makecell[c]{0.071\\{(0.047)}} & \makecell[c]{0.27***\\{(0.055)}} & \makecell[c]{0.199***\\{(0.069)}}\\
Calendar & \makecell[c]{0.052\\{(0.042)}} & \makecell[c]{0.037\\{(0.055)}} & \makecell[c]{0.071\\{(0.059)}} & \makecell[c]{0.033\\{(0.078)}}\\
Ink & \makecell[c]{0.112***\\{(0.04)}} & \makecell[c]{0.063\\{(0.051)}} & \makecell[c]{0.175***\\{(0.056)}} & \makecell[c]{0.112\\{(0.072)}}\\
\midrule
Control & \makecell[c]{0.671\\{(0.031)}} & \makecell[c]{0.735\\{(0.042)}} & \makecell[c]{0.59\\{(0.047)}} & \makecell[c]{-0.144**\\{(0.063)}}\\
Observations & 1,141 & 1,141 & 1,141 & 1,141\\
$H0$: Any Signal = No Signal, $p$-value & $<$0.001 & 0.145 & $<$0.001 & 0.002\\
$H0$: Bracelet = Calendar, $p$-value & $<$0.001 & 0.389 & $<$0.001 & 0.003\\
\bottomrule
\end{tabular}}


	}
    \floatfoot{\textit{Notes:}
    This table replicates Table~\ref{tab:beliefs-visibility}, replacing the binary Close/Far indicator with continuous community-centroid-to-PoT distance and its interactions with item/signal arm indicators. The dependent variable is the respondent-level observability measure: the fraction of recognized community members whose deworming status the respondent reports knowing. Rows for Bracelet, Calendar, and Ink report effects relative to Control. The Control row reports the implied Control mean in columns (1)--(3) and the implied Control Far--Close difference in column (4). Columns labeled Close and Far report implied treatment effects averaged over observations in Close and Far communities. Under the linear continuous distance specification, this is equivalent to evaluating effects at the mean community-centroid-to-PoT distance in each cell. Column (4) reports the implied Far--Close difference. We pool Bracelet and Ink as ``Any Signal'' and Control and Calendar as ``No Signal'' to test \(H_0\): Any Signal \(=\) No Signal. The second test, \(H_0\): Bracelet \(=\) Calendar, compares the bracelet to the active comparison item. All regressions include county fixed effects, LASSO-selected controls for sex and age, and expected community-centroid-to-PoT distance as a design control. Bootstrapped standard errors, clustered at the community level, are reported in parentheses. \(\sym{*}p<0.10\), \(\sym{**}p<0.05\), \(\sym{***}p<0.01\).}
\end{threeparttable}
\end{table}

\begin{table}[ht!]
	\scriptsize
	\caption{Effects of Item/Signal Arms and Distance on the Observability of Deworming: Without Individual or Expected Distance Controls}
    \label{tab:rf-fob-no-covs-no-mud}
	\centering
    \begin{threeparttable}
    \resizebox{\textwidth}{!}{%
    
\centering

\centering
\resizebox{\ifdim\width>\linewidth\linewidth\else\width\fi}{!}{
\begin{tabular}[t]{lcccc}
\toprule
\multicolumn{1}{c}{ } & \multicolumn{4}{c}{Reduced Form} \\
\cmidrule(l{3pt}r{3pt}){2-5}
\multicolumn{1}{c}{ } & \multicolumn{1}{c}{Combined} & \multicolumn{1}{c}{Close} & \multicolumn{1}{c}{Far} & \multicolumn{1}{c}{Far - Close} \\
Dependent variable: Observability & (1) & (2) & (3) & (4)\\
\midrule
Bracelet & \makecell[c]{0.143***\\{(0.038)}} & \makecell[c]{0.063\\{(0.048)}} & \makecell[c]{0.246***\\{(0.057)}} & \makecell[c]{0.183**\\{(0.073)}}\\
Calendar & \makecell[c]{0.04\\{(0.043)}} & \makecell[c]{0.02\\{(0.058)}} & \makecell[c]{0.066\\{(0.063)}} & \makecell[c]{0.046\\{(0.087)}}\\
Ink & \makecell[c]{0.1**\\{(0.04)}} & \makecell[c]{0.045\\{(0.053)}} & \makecell[c]{0.169***\\{(0.059)}} & \makecell[c]{0.124\\{(0.078)}}\\
\midrule
Control & \makecell[c]{0.683\\{(0.032)}} & \makecell[c]{0.749\\{(0.041)}} & \makecell[c]{0.598\\{(0.05)}} & \makecell[c]{-0.15**\\{(0.063)}}\\
Observations & 1,141 & 1,141 & 1,141 & 1,141\\
$H0$: Any Signal = No Signal, $p$-value & $<$0.001 & 0.221 & $<$0.001 & 0.01\\
$H0$: Bracelet = Calendar, $p$-value & 0.001 & 0.352 & $<$0.001 & 0.043\\
\bottomrule
\end{tabular}}


	}
    \floatfoot{\textit{Notes:}
    This table replicates Table~\ref{tab:beliefs-visibility}, omitting the LASSO-selected controls for sex and age and the expected community-centroid-to-PoT distance design control. The specification retains county fixed effects. Bootstrapped standard errors, clustered at the community level, are reported in parentheses. \(\sym{*}p<0.10\), \(\sym{**}p<0.05\), \(\sym{***}p<0.01\).
}
\end{threeparttable}
\end{table}




\begin{table}[ht!]
\scriptsize
\caption{Effects of Item/Signal Arms and Distance on the Observability of Deworming, Controlling for Unbalanced Baseline Cluster Covariates}
\label{tab:fob-baseline-imbalance-robustness}
\centering
\begin{threeparttable}
\begin{tabular}{lcccc}
\toprule
 & Combined & Close & Far & Far--Close \\
\multicolumn{5}{l}{\textbf{First-order beliefs}} \\
\addlinespace[0.2em]
Bracelet & \makecell[c]{0.147***\\{(0.039)}} & \makecell[c]{0.070\\{(0.049)}} & \makecell[c]{0.244***\\{(0.061)}} & \makecell[c]{0.182**\\{(0.080)}} \\
Calendar & \makecell[c]{0.041\\{(0.046)}} & \makecell[c]{0.037\\{(0.063)}} & \makecell[c]{0.051\\{(0.064)}} & \makecell[c]{0.028\\{(0.091)}} \\
Ink & \makecell[c]{0.104**\\{(0.044)}} & \makecell[c]{0.058\\{(0.050)}} & \makecell[c]{0.162**\\{(0.070)}} & \makecell[c]{0.113\\{(0.084)}} \\
Control & \makecell[c]{0.684\\{(0.038)}} & \makecell[c]{0.749\\{(0.048)}} & \makecell[c]{0.599\\{(0.049)}} & \makecell[c]{-0.133*\\{(0.068)}} \\
\addlinespace[0.3em]
Baseline knows medicine treats worms share & \makecell[c]{-0.058\\{(0.105)}} & \makecell[c]{-0.090\\{(0.158)}} & \makecell[c]{0.069\\{(0.126)}} & \makecell[c]{-0.003\\{(0.098)}} \\
Baseline past-year dewormed share & \makecell[c]{0.056\\{(0.087)}} & \makecell[c]{0.003\\{(0.111)}} & \makecell[c]{0.069\\{(0.135)}} & \makecell[c]{0.017\\{(0.088)}} \\
Baseline knows bi-yearly treatment share & \makecell[c]{0.066\\{(0.107)}} & \makecell[c]{0.078\\{(0.145)}} & \makecell[c]{0.029\\{(0.183)}} & \makecell[c]{0.047\\{(0.112)}} \\
\addlinespace[0.3em]
Observations & 1,141 & 641 & 500 & 1,141 \\
Clusters & 144 & 80 & 64 & 144 \\
County fixed effects & Yes & Yes & Yes & Yes \\
Main covariates & Yes & Yes & Yes & Yes \\
Unbalanced baseline covariates & Yes & Yes & Yes & Yes \\
\bottomrule
\end{tabular}

\floatfoot{\textit{Notes:}
This table replicates Table~\ref{tab:beliefs-visibility} while adding controls for the unbalanced baseline survey covariates, discussed in Section \ref{app:baseline-sample}. Because baseline respondents cannot be linked to individuals in the endline observability sample, these covariates are measured as community-level shares among baseline survey respondents: knowledge that medicine treats worms, deworming in the past year, and knowledge that biannual treatment is recommended. All regressions include county fixed effects, LASSO-selected controls for sex and age, and expected community-centroid-to-PoT distance as a design control. Bootstrapped standard errors, clustered at the community level, are reported in parentheses. \(\sym{*}p<0.10\), \(\sym{**}p<0.05\), \(\sym{***}p<0.01\).}
\end{threeparttable}
\end{table}


\begin{table}[ht!]
	\scriptsize
	\caption{Reported Knowledge and Accuracy of Peer Deworming Knowledge}
    \label{tab:rf-fob-decomp}
	\centering
    \begin{threeparttable}
    \resizebox{\textwidth}{!}{%
    \centering
\resizebox{\ifdim\width>\linewidth\linewidth\else\width\fi}{!}{
\begin{tabular}[t]{lcccc}
\toprule
\multicolumn{1}{c}{ } & \multicolumn{4}{c}{Reduced Form} \\
\cmidrule(l{3pt}r{3pt}){2-5}
\multicolumn{1}{c}{ } 
& \multicolumn{1}{c}{Combined} 
& \multicolumn{1}{c}{Close} 
& \multicolumn{1}{c}{Far} 
& \multicolumn{1}{c}{Far--Close} \\
\midrule

\multicolumn{5}{l}{\textit{Panel A: Reported knowledge (definite Yes/No answer)}} \\
Bracelet & \makecell[c]{0.149***\\{(0.038)}} & \makecell[c]{0.066\\{(0.049)}} & \makecell[c]{0.256***\\{(0.056)}} & \makecell[c]{0.19***\\{(0.073)}}\\
Calendar & \makecell[c]{0.044\\{(0.043)}} & \makecell[c]{0.027\\{(0.06)}} & \makecell[c]{0.065\\{(0.062)}} & \makecell[c]{0.038\\{(0.087)}}\\
Ink & \makecell[c]{0.103**\\{(0.04)}} & \makecell[c]{0.053\\{(0.054)}} & \makecell[c]{0.167***\\{(0.06)}} & \makecell[c]{0.114\\{(0.08)}}\\
\midrule
Control mean & \makecell[c]{0.679\\{(0.033)}} & \makecell[c]{0.744\\{(0.043)}} & \makecell[c]{0.596\\{(0.049)}} & \makecell[c]{-0.148**\\{(0.063)}}\\
Observations & 1,141 & 1,141 & 1,141 & 1,141\\
\(H_0\): Any Signal = No Signal, \(p\) value & $<$0.001 & 0.2 & $<$0.001 & 0.008\\
\(H_0\): Bracelet = Calendar, \(p\) value & $<$0.001 & 0.403 & $<$0.001 & 0.027\\
\midrule

\multicolumn{5}{l}{\textit{Panel B: Accuracy conditional on definite Yes/No answer}} \\
Bracelet & \makecell[c]{-0.043\\{(0.027)}} & \makecell[c]{-0.007\\{(0.039)}} & \makecell[c]{-0.09**\\{(0.038)}} & \makecell[c]{-0.083\\{(0.055)}}\\
Calendar & \makecell[c]{-0.048\\{(0.031)}} & \makecell[c]{-0.023\\{(0.04)}} & \makecell[c]{-0.081*\\{(0.049)}} & \makecell[c]{-0.058\\{(0.064)}}\\
Ink & \makecell[c]{-0.088**\\{(0.035)}} & \makecell[c]{-0.102**\\{(0.049)}} & \makecell[c]{-0.07\\{(0.049)}} & \makecell[c]{0.032\\{(0.069)}}\\
\midrule
Control mean & \makecell[c]{0.656\\{(0.019)}} & \makecell[c]{0.65\\{(0.024)}} & \makecell[c]{0.664\\{(0.031)}} & \makecell[c]{0.013\\{(0.04)}}\\
Observations & 1,035 & 1,035 & 1,035 & 1,035\\
\(H_0\): Any Signal = No Signal, \(p\) value & 0.078 & 0.199 & 0.212 & 0.937\\
\(H_0\): Bracelet = Calendar, \(p\) value & 0.862 & 0.714 & 0.829 & 0.689\\
\midrule

\multicolumn{5}{l}{\textit{Panel C: Correct classification (``Don't know'' counted as incorrect)}} \\
Bracelet & \makecell[c]{0.069**\\{(0.032)}} & \makecell[c]{0.029\\{(0.044)}} & \makecell[c]{0.119***\\{(0.042)}} & \makecell[c]{0.091\\{(0.059)}}\\
Calendar & \makecell[c]{-0.003\\{(0.032)}} & \makecell[c]{-0.006\\{(0.046)}} & \makecell[c]{0.001\\{(0.043)}} & \makecell[c]{0.007\\{(0.063)}}\\
Ink & \makecell[c]{-0.012\\{(0.031)}} & \makecell[c]{-0.064\\{(0.043)}} & \makecell[c]{0.053\\{(0.044)}} & \makecell[c]{0.118*\\{(0.062)}}\\
\midrule
Control mean & \makecell[c]{0.437\\{(0.023)}} & \makecell[c]{0.483\\{(0.032)}} & \makecell[c]{0.378\\{(0.032)}} & \makecell[c]{-0.105**\\{(0.045)}}\\
Observations & 1,134 & 1,134 & 1,134 & 1,134\\
\(H_0\): Any Signal = No Signal, \(p\) value & 0.169 & 0.628 & 0.003 & 0.015\\
\(H_0\): Bracelet = Calendar, \(p\) value & 0.011 & 0.401 & 0.003 & 0.162\\
\bottomrule
\end{tabular}}
	}
    \floatfoot{\textit{Notes:}
    Each panel reports OLS estimates using the named-individual peer belief items in the no-SMS endline observability sample. Panel A reports the definite-answer rate: the share of recognized peers for whom the respondent gives a Yes or No response (1,141 respondents; administrative take-up need not be observed). Panel B reports conditional accuracy: the share of Yes or No responses with observed administrative take-up that are correct (1,035 respondents). It excludes 91 respondents with no definite answers and 15 whose definite answers all lack administrative take-up records. Panel C reports correct observability, counting ``Don't know'' as incorrect even when administrative take-up is missing and omitting Yes or No responses without administrative take-up records (1,134 respondents). It excludes seven respondents with neither a checkable definite answer nor a ``Don't know'' response.  The first hypothesis test pools Bracelet and Ink as ``Any Signal'' and Control and Calendar as ``No Signal'' to test \(H_0\): Any Signal \(=\) No Signal. The second test compares Bracelet to Calendar. All regressions include county fixed effects, LASSO selected controls for sex and age, and expected community-centroid-to-PoT distance as a design control. Bootstrapped standard errors, clustered at the community level, are reported in parentheses.  \(^{*}p<0.10\), \(^{**}p<0.05\), \(^{***}p<0.01\).}
\end{threeparttable}
\end{table}


\begin{table}[ht!]
	\scriptsize
	\caption{Effects of Item/Signal Arms and Distance on the Perceived Observability of Deworming}
    \label{tab:rf-sob-ates}
	\centering
    \begin{threeparttable}
    \resizebox{\textwidth}{!}{%
    
\centering

\centering
\resizebox{\ifdim\width>\linewidth\linewidth\else\width\fi}{!}{
\begin{tabular}[t]{lcccc}
\toprule
\multicolumn{1}{c}{ } & \multicolumn{4}{c}{Reduced Form} \\
\cmidrule(l{3pt}r{3pt}){2-5}
\multicolumn{1}{c}{ } & \multicolumn{1}{c}{Combined} & \multicolumn{1}{c}{Close} & \multicolumn{1}{c}{Far} & \multicolumn{1}{c}{Far - Close} \\
Dependent variable: Observability Beliefs & (1) & (2) & (3) & (4)\\
\midrule
Bracelet & \makecell[c]{0.088*\\{(0.053)}} & \makecell[c]{0.066\\{(0.067)}} & \makecell[c]{0.116\\{(0.079)}} & \makecell[c]{0.05\\{(0.1)}}\\
Calendar & \makecell[c]{0.029\\{(0.052)}} & \makecell[c]{0.054\\{(0.072)}} & \makecell[c]{-0.004\\{(0.075)}} & \makecell[c]{-0.058\\{(0.103)}}\\
Ink & \makecell[c]{0.06\\{(0.052)}} & \makecell[c]{0.031\\{(0.072)}} & \makecell[c]{0.097\\{(0.074)}} & \makecell[c]{0.066\\{(0.101)}}\\
\midrule
Control & \makecell[c]{0.634\\{(0.042)}} & \makecell[c]{0.672\\{(0.055)}} & \makecell[c]{0.585\\{(0.059)}} & \makecell[c]{-0.088\\{(0.078)}}\\
Observations & 1,141 & 1,141 & 1,141 & 1,141\\
$H0$: Any Signal = No Signal, $p$-value & 0.084 & 0.635 & 0.037 & 0.204\\
$H0$: Bracelet = Calendar, $p$-value & 0.142 & 0.833 & 0.08 & 0.256\\
\bottomrule
\end{tabular}}


	}
     \floatfoot{\textit{Notes:} This table replicates the specification in Table~\ref{tab:beliefs-visibility} using the respondent level measure of perceived observability: the fraction of recognized community members whom the respondent reports know the respondent's own deworming status. \(\sym{*}p<0.10\), \(\sym{**}p<0.05\), \(\sym{***}p<0.01\).}
\end{threeparttable}
\end{table}


\begin{table}[ht!]
	\scriptsize
	\caption{Endline Predictions about Community Deworming Take-up}
    \label{tab:pred-te}
	\centering
    \begin{threeparttable}
    \resizebox{\textwidth}{!}{%
    \centering
\resizebox{\ifdim\width>\linewidth\linewidth\else\width\fi}{!}{
\begin{tabular}[t]{lcccc}
\toprule
\multicolumn{1}{c}{ } & \multicolumn{4}{c}{Reduced Form} \\
\cmidrule(l{3pt}r{3pt}){2-5}
\multicolumn{1}{c}{ } & \multicolumn{1}{c}{Combined} & \multicolumn{1}{c}{Close} & \multicolumn{1}{c}{Far} & \multicolumn{1}{c}{Far - Close} \\
Dependent variable: Predicted Take-up & (1) & (2) & (3) & (4)\\
\midrule
Bracelet & \makecell[c]{0.104***\\{(0.023)}} & \makecell[c]{0.065**\\{(0.026)}} & \makecell[c]{0.152***\\{(0.04)}} & \makecell[c]{0.087*\\{(0.048)}}\\
Calendar & \makecell[c]{0.058**\\{(0.023)}} & \makecell[c]{0.024\\{(0.029)}} & \makecell[c]{0.101**\\{(0.04)}} & \makecell[c]{0.077\\{(0.051)}}\\
Ink & \makecell[c]{0.029\\{(0.024)}} & \makecell[c]{0.009\\{(0.025)}} & \makecell[c]{0.055\\{(0.045)}} & \makecell[c]{0.047\\{(0.053)}}\\
\midrule
Control & \makecell[c]{0.559\\{(0.019)}} & \makecell[c]{0.617\\{(0.019)}} & \makecell[c]{0.488\\{(0.038)}} & \makecell[c]{-0.129***\\{(0.044)}}\\
Observations & 2,659 & 2,659 & 2,659 & 2,659\\
\bottomrule
\end{tabular}}


	}
    \floatfoot{\textit{Notes:}
    Each column reports OLS estimates using respondents' endline predictions of community deworming take-up. At endline, respondents were asked: ``Out of 10 people in your community, how many do you think came for deworming?'' The dependent variable is the stated number divided by 10. The estimating equation is the same as in Table~\ref{tab:overall_effects}. The Control row reports the Control mean in columns (1)--(3) and the Control Far--Close difference in column (4). All regressions include county fixed effects, LASSO selected controls for sex and age, and expected community-centroid-to-PoT distance as a design control. Bootstrapped standard errors, clustered at the community level, are reported in parentheses. \(\sym{*}p<0.10\), \(\sym{**}p<0.05\), \(\sym{***}p<0.01\).}
    \end{threeparttable}
    \end{table}


\begin{table}[ht!]
	\scriptsize
	\caption{Effects of Item/Signal Arms and Continuous Distance on Deworming Take-up}
    \label{tab:rf-ates-cts-dist}
	\centering
    \begin{threeparttable}
    \resizebox{\textwidth}{!}{%
    \centering
\resizebox{\ifdim\width>\linewidth\linewidth\else\width\fi}{!}{
\begin{tabular}[t]{lcccc}
\toprule
\multicolumn{1}{c}{ } & \multicolumn{4}{c}{Reduced Form} \\
\cmidrule(l{3pt}r{3pt}){2-5}
\multicolumn{1}{c}{ } & \multicolumn{1}{c}{Combined} & \multicolumn{1}{c}{Close} & \multicolumn{1}{c}{Far} & \multicolumn{1}{c}{Far - Close} \\
Dependent variable: Take-up & (1) & (2) & (3) & (4)\\
\midrule
Bracelet & \makecell[c]{0.087***\\{(0.028)}} & \makecell[c]{0.062*\\{(0.033)}} & \makecell[c]{0.118***\\{(0.045)}} & \makecell[c]{0.056\\{(0.053)}}\\
Calendar & \makecell[c]{0.037\\{(0.027)}} & \makecell[c]{0.029\\{(0.03)}} & \makecell[c]{0.048\\{(0.037)}} & \makecell[c]{0.02\\{(0.038)}}\\
Ink & \makecell[c]{-0.013\\{(0.031)}} & \makecell[c]{-0.049\\{(0.031)}} & \makecell[c]{0.032\\{(0.042)}} & \makecell[c]{0.081**\\{(0.041)}}\\
\midrule
Control & \makecell[c]{0.327\\{(0.023)}} & \makecell[c]{0.405\\{(0.023)}} & \makecell[c]{0.229\\{(0.033)}} & \makecell[c]{-0.176***\\{(0.032)}}\\
Observations & 9,805 & 9,805 & 9,805 & 9,805\\
$H0$: Any Signal $=$ No Signal, $p$-value & 0.317 & 0.712 & 0.061 & 0.049\\
$H0$: Bracelet $=$ Calendar, $p$-value & 0.032 & 0.275 & 0.034 & 0.407\\
\bottomrule
\end{tabular}}


	}
    \floatfoot{\textit{Notes:}
    This table replicates Table~\ref{tab:overall_effects}, replacing the binary Close/Far indicator with continuous community-centroid-to-PoT distance and its interactions with item/signal arm indicators. The dependent variable is an indicator equal to one if the adult dewormed during the campaign. Rows for Bracelet, Calendar, and Ink report effects relative to Control. The Control row reports the implied Control mean in columns (1)--(3) and the implied Control Far--Close difference in column (4). Columns labeled Close and Far report implied treatment effects averaged over observations in Close and Far communities. Under the linear continuous-distance specification, this is equivalent to evaluating effects at the mean community-centroid-to-PoT distance in each cell. Column (4) reports the implied Far--Close difference. We pool Bracelet and Ink as ``Any Signal'' and Control and Calendar as ``No Signal'' to test \(H_0\): Any Signal \(=\) No Signal. The second test, \(H_0\): Bracelet \(=\) Calendar, compares the bracelet to the active comparison item. All regressions include county fixed effects, LASSO-selected controls for sex and age, and expected community-centroid-to-PoT distance as a design control. Bootstrapped standard errors, clustered at the community level, are reported in parentheses. \(\sym{*}p<0.10\), \(\sym{**}p<0.05\), \(\sym{***}p<0.01\).
}
\end{threeparttable}
\end{table}

\begin{table}[ht!]
	\scriptsize
	\caption{Effects of Item/Signal Arms and Distance on Deworming Take-up: Without Individual or Expected Distance Controls}
    \label{tab:rf-takeup-no-covs-no-mud}
	\centering
    \begin{threeparttable}
    \resizebox{\textwidth}{!}{%
    
\centering

\centering
\resizebox{\ifdim\width>\linewidth\linewidth\else\width\fi}{!}{
\begin{tabular}[t]{lcccc}
\toprule
\multicolumn{1}{c}{ } & \multicolumn{4}{c}{Reduced Form} \\
\cmidrule(l{3pt}r{3pt}){2-5}
\multicolumn{1}{c}{ } & \multicolumn{1}{c}{Combined} & \multicolumn{1}{c}{Close} & \multicolumn{1}{c}{Far} & \multicolumn{1}{c}{Far - Close} \\
Dependent variable: Take-up & (1) & (2) & (3) & (4)\\
\midrule
Bracelet & \makecell[c]{0.074***\\{(0.029)}} & \makecell[c]{0.041\\{(0.038)}} & \makecell[c]{0.116**\\{(0.046)}} & \makecell[c]{0.074\\{(0.06)}}\\
Calendar & \makecell[c]{0.022\\{(0.029)}} & \makecell[c]{0.02\\{(0.039)}} & \makecell[c]{0.024\\{(0.041)}} & \makecell[c]{0.004\\{(0.055)}}\\
Ink & \makecell[c]{-0.034\\{(0.032)}} & \makecell[c]{-0.061\\{(0.039)}} & \makecell[c]{-0.001\\{(0.052)}} & \makecell[c]{0.059\\{(0.063)}}\\
\midrule
Control & \makecell[c]{0.34\\{(0.024)}} & \makecell[c]{0.414\\{(0.028)}} & \makecell[c]{0.25\\{(0.038)}} & \makecell[c]{-0.164***\\{(0.046)}}\\
Observations & 9,805 & 9,805 & 9,805 & 9,805\\
$H0$: Any Signal = No Signal, $p$-value & 0.64 & 0.431 & 0.141 & 0.102\\
$H0$: Bracelet = Calendar, $p$-value & 0.032 & 0.555 & 0.003 & 0.136\\
\bottomrule
\end{tabular}}


	}
    \floatfoot{\textit{Notes:}
This table replicates Table~\ref{tab:overall_effects}, omitting the LASSO-selected controls for sex and age and the expected community-centroid-to-PoT distance design control. The specification retains county fixed effects. Bootstrapped standard errors, clustered at the community level, are reported in parentheses. \(\sym{*}p<0.10\), \(\sym{**}p<0.05\), \(\sym{***}p<0.01\).
}
\end{threeparttable}
\end{table}




\begin{table}[ht!]
\scriptsize
\caption{Effects of Item/Signal Arms and Distance on Deworming Takeup, Controlling for Unbalanced Baseline Cluster Covariates}
\label{tab:prior-deworming-robustness}
\centering
\begin{threeparttable}
\input{tables/prior-deworming-robustness}
\floatfoot{\textit{Notes:}
This table replicates Table~\ref{tab:overall_effects} while adding controls for the unbalanced baseline survey covariates, discussed in Section \ref{app:baseline-sample}. Because baseline respondents cannot be linked to individuals in the take-up sample, these covariates are measured as community-level shares among baseline survey respondents: knowledge that medicine treats worms, deworming in the past year, and knowledge that biannual treatment is recommended. All regressions include county fixed effects, LASSO-selected controls for sex and age, and expected community-centroid-to-PoT distance as a design control. Bootstrapped standard errors, clustered at the community level, are reported in parentheses. \(\sym{*}p<0.10\), \(\sym{**}p<0.05\), \(\sym{***}p<0.01\).}
\end{threeparttable}
\end{table}



%\floatfoot{\textit{Notes:} Each column reports heterogeneous treatment effects for a single covariate group. We estimate $Y_{ic}=X_{ic}\beta+\sum_{z=1}^{4}\bigl(\gamma_z\,\text{treatment}_{icz}\times X_{ic} +\beta_z\,\text{treatment}_{icz}\bigr)+\alpha_s$, where \(Y_{ic}\) is a deworming indicator, \(X_{ic}\) is the covariate listed in the column header, and \(\alpha_s\) is a county‑strata fixed effect. Age \(>40\) – individual is older than 40. Female – individual is female. Phone Owner – individual owns a phone.  Community Judgment – community’s praise/stigma score above the sample mean. Previously Dewormed – community’s prior deworming rate above the sample mean. Externality Knowledge – community’s knowledge of externalities score above the mean. Standard errors are clustered at the community level. \sym{*}$p<0.10$, \sym{**}$p<0.05$, \sym{***}$p<0.01$.}

\begin{table}[ht!]
	\small
	\caption{Treatment Effect Heterogeneity by Baseline Covariates}
    \label{tab:rf-covs}
	\centering
    \begin{threeparttable}
    \resizebox{\textwidth}{!}{%
    \input{rf-tables/het-tes}
	}
\floatfoot{\textit{Notes:}
Each column reports a linear probability model for deworming take-up. The covariate \(X\) interacted with treatment assignment is listed in the table header. Rows for Bracelet, Calendar, and Ink report treatment effects relative to Control for observations with \(X=0\). The row \(X\) reports the association between the covariate and take-up in the Control arm. Interaction rows report how each treatment effect differs for observations with \(X=1\). All regressions include county fixed effects. Age \(>40\), Female, and Phone owner are individual-level indicators. High social judgment indicates that the community's baseline praise/stigma score is above the sample mean. High prior deworming indicates that the community's prior deworming rate is above the sample mean. High externality knowledge indicates that the community's baseline knowledge of worm transmission externalities is above the sample mean. Standard errors are clustered at the community level. \(\sym{*}p<0.10\), \(\sym{**}p<0.05\), \(\sym{***}p<0.01\).}
\end{threeparttable}
\end{table}


\begin{table}[ht!]
	\small
    \caption{Effects of Reminder and Social Information SMS on Take-up in Control Communities}
    \label{tab:sms-reminder-te}
	\centering
    \begin{threeparttable}
    % \resizebox{\textwidth}{!}{%
    \input{rf-tables/main-specs/incentive-control-sms-te}
	% }
    \floatfoot{\textit{Notes:}
    This table reports estimates from a linear probability model (OLS) with a binary outcome for deworming take‑up.  The regressors are indicators for receiving a Social Information SMS or a Reminder Only SMS. The omitted group is phone owners who received no SMS. Because the Reminder treatment was implemented only in Control communities, the analysis sample is restricted to phone owners in the Control arm, not the full SMS treatment sample ($N$=3{,}022).  All regressions include county fixed effects and the LASSO selected controls for sex and age, and expected distance as a design control. Bootstrapped standard errors, clustered at the community level, are in parentheses.\(\sym{*}p<0.10\), \(\sym{**}p<0.05\), \(\sym{***}p<0.01\).}
\end{threeparttable}
\end{table}

\begin{table}[ht!]
	\small
	\caption{Externality Knowledge at Endline}
    \label{tab:endline-deworm-knowledge}
	\centering
    \begin{threeparttable}
    
\centering

\centering
\resizebox{\ifdim\width>\linewidth\linewidth\else\width\fi}{!}{
\begin{tabular}[t]{lcccc}
\toprule
\multicolumn{1}{c}{ } & \multicolumn{4}{c}{Reduced Form} \\
\cmidrule(l{3pt}r{3pt}){2-5}
\multicolumn{1}{c}{ } & \multicolumn{1}{c}{Combined} & \multicolumn{1}{c}{Close} & \multicolumn{1}{c}{Far} & \multicolumn{1}{c}{Far - Close} \\
Dependent variable: Externality Knowledge & (1) & (2) & (3) & (4)\\
\midrule
Bracelet & \makecell[c]{0.033\\{(0.036)}} & \makecell[c]{0.002\\{(0.05)}} & \makecell[c]{0.071\\{(0.052)}} & \makecell[c]{0.069\\{(0.072)}}\\
Calendar & \makecell[c]{0.032\\{(0.036)}} & \makecell[c]{0.038\\{(0.054)}} & \makecell[c]{0.025\\{(0.045)}} & \makecell[c]{-0.014\\{(0.069)}}\\
Ink & \makecell[c]{0.095**\\{(0.039)}} & \makecell[c]{0.112**\\{(0.051)}} & \makecell[c]{0.074\\{(0.058)}} & \makecell[c]{-0.038\\{(0.077)}}\\
\midrule
Control & \makecell[c]{0.3\\{(0.027)}} & \makecell[c]{0.301\\{(0.042)}} & \makecell[c]{0.299\\{(0.032)}} & \makecell[c]{-0.002\\{(0.053)}}\\
Observations & 2,659 & 2,659 & 2,659 & 2,659\\
$H0$: Any Signal = No Signal, $p$-value & 0.061 & 0.266 & 0.118 & 0.662\\
$H0$: Bracelet = Calendar, $p$-value & 0.978 & 0.419 & 0.351 & 0.228\\
\bottomrule
\end{tabular}}


\floatfoot{%
\textit{Notes:} Each column reports OLS estimates of the respondent-level externality knowledge measure, equal to 1 if the respondent answers `Yes’ to `Can a person sick with worms spread worms to others?' or to both of the following: `If you have worms, does that affect your neighbors’ or relatives’ health?' and `If your neighbors or relatives have worms, does that affect your health?'. “Control’’ is the mean in the Control arm.  Rows for Bracelet, Ink, and Calendar give treatment effects relative to Control.  Column (4) reports the difference in those effects between Far and Close communities.  We pool Bracelet and Ink as ``Any Signal'' and Control and Calendar as ``No Signal'' to test $H_0$: Any Signal $=$ No Signal.  The second test, $H_0$: Bracelet$=$Calendar, compares the bracelet and the active comparison item. All regressions include county fixed effects and the LASSO selected controls for sex and age, and expected distance to the PoT. Bootstrapped standard errors, clustered at the community level, are reported in parentheses.\sym{*}$p<0.10$, \sym{**}$p<0.05$, \sym{***}$p<0.01$.}
\end{threeparttable}
\end{table}


\begin{table}[ht!]
  \small
  \caption{Difference in Monetary Valuation of Calendar vs.\ Bracelet}
  \label{tab:wtp-summ-table}
  \centering
  \begin{threeparttable}
      % \resizebox{\textwidth}{!}{%
    \input{tables/wtp-summ-table}  % ← underlying values unchanged
    % }
    \floatfoot{\textit{Notes:}
    The table reports posterior summaries for the relative valuation parameter \(\psi\) from the joint structural estimation. The parameter \(\psi\) measures the mean calendar-minus-bracelet valuation difference in US dollars, the scaled monetary units used in the elicitation likelihood. Positive values indicate that respondents privately value the calendar more than the bracelet. Panel B reports posterior predicted probabilities that a respondent prefers the calendar across multiple hypothetical switch offers. The two-stage probit likelihood is detailed in Appendix~\ref{sec:wtp-model}. In the take-up model, \(\psi\) is linked to the calendar and bracelet private payoff coefficients through the restriction \(\beta_{\text{Calendar}}=\beta_{\text{Bracelet}}+\rho\psi\). The joint model combines observations from the take-up analysis sample (\(N=9{,}805\)), the relative valuation sample (\(N=467\)), and the endline observability sample (\(N=1{,}141\)).}
      %Survey wording: respondents first chose between the items (``Which one would you like— the calendar or the bracelet?’’), then faced a randomly drawn cash offer \(x\) (``You can take the [bracelet/calendar] on its own, \emph{or} the [calendar/bracelet] and \(x\) KSh’’).}
  \end{threeparttable}
\end{table}


\begin{table}[ht!]
	\small
	\caption{Gift‑Choice Preference for a Bracelet vs. Calendar by Distance}
    \label{tab:pref-cal-gift}
	\centering
    \begin{threeparttable}
    \resizebox{\textwidth}{!}{%
    
\centering

\centering
\resizebox{\ifdim\width>\linewidth\linewidth\else\width\fi}{!}{
\begin{tabular}[t]{lcccc}
\toprule
\multicolumn{1}{c}{ } & \multicolumn{4}{c}{Reduced Form} \\
\cmidrule(l{3pt}r{3pt}){2-5}
\multicolumn{1}{c}{ } & \multicolumn{1}{c}{Combined} & \multicolumn{1}{c}{Close} & \multicolumn{1}{c}{Far} & \multicolumn{1}{c}{Far - Close} \\
Dependent variable: Prefer Bracelet & (1) & (2) & (3) & (4)\\
\midrule
Bracelet & \makecell[c]{-0.093***\\{(0.032)}} & \makecell[c]{-0.102**\\{(0.04)}} & \makecell[c]{-0.086*\\{(0.047)}} & \makecell[c]{0.016\\{(0.062)}}\\
Calendar & \makecell[c]{0.174***\\{(0.043)}} & \makecell[c]{0.22***\\{(0.061)}} & \makecell[c]{0.131**\\{(0.057)}} & \makecell[c]{-0.089\\{(0.081)}}\\
Ink & \makecell[c]{0.029\\{(0.037)}} & \makecell[c]{0.048\\{(0.045)}} & \makecell[c]{0.01\\{(0.056)}} & \makecell[c]{-0.038\\{(0.071)}}\\
\midrule
Control & \makecell[c]{0.227\\{(0.021)}} & \makecell[c]{0.221\\{(0.028)}} & \makecell[c]{0.232\\{(0.031)}} & \makecell[c]{0.011\\{(0.041)}}\\
$|\text{Calendar}| - |\text{Bracelet}|$ & \makecell[c]{0.08\\{(0.061)}} & \makecell[c]{0.118\\{(0.083)}} & \makecell[c]{0.045\\{(0.081)}} & \makecell[c]{-0.073\\{(0.113)}}\\
Observations & 1,174 & 1,174 & 1,174 & 1,174\\
\bottomrule
\end{tabular}}


	}
    \floatfoot{\textit{Notes:} The table examines whether preferences between the bracelet and calendar vary across item/signal arms and with distance (Far vs.\ Close). The dependent variable equals 1 if a respondent chooses the bracelet rather than the calendar when offered a choice between the two items as an endline gift. The sample is drawn from the endline sample (\(N=2{,}659\)) and restricted to adults who did not deworm and therefore did not receive either item during the campaign (\(N=1{,}174\)). In Control communities, where neither item was distributed during the campaign, gift choices are informative about relative private consumption value among non-dewormers. Comparisons across item/signal arms may additionally reflect post-treatment selection and any utility or disutility associated with local item prevalence and the local signaling environment. Estimates follow the main reduced form specification: OLS with strata fixed effects, LASSO-selected controls (sex, age, and expected distance to the point of treatment), and standard errors bootstrapped and clustered at the community level. Rows for Bracelet, Ink, and Calendar report effects relative to Control in each column; Column (4) reports the Far--Close difference in those effects. \( |\text{Calendar}|-|\text{Bracelet}| \) gives the difference in the magnitudes of the Calendar and Bracelet coefficients. In Column (4), this row reports how that difference changes between Close and Far communities.\sym{*}$p<0.10$, \sym{**}$p<0.05$, \sym{***}$p<0.01$.}
\end{threeparttable}
\end{table}


\clearpage

\newpage

\begin{table}[ht!]
  \scriptsize
  \caption{Posterior Sensitivity to Likelihood Reweighting}
  \label{tab:struct-sensitivity}
  \centering
  \begin{threeparttable}
    \input{tables/structural-object-sensitivity-1500m}
\floatfoot{\textit{Notes:}
This table reports Bayesian posterior sensitivity to local reweighting of the
  take-up, observability, and WTP components of the likelihood.   For structural object \(Q(\theta)\) and likelihood component \(b\), each entry
  is computed from the joint posterior draws as
  \(\left|\operatorname{Cov}(Q(\theta),\ell_b(\theta)\mid Y)\right|/
  \operatorname{sd}(Q(\theta)\mid Y)\), where \(\ell_b(\theta)\) is the block’s
  log likelihood. This equals the absolute local change in the posterior mean,
  in posterior-standard-deviation units, when the block’s likelihood weight is
  marginally increased from one. Entries are
  absolute standardized changes in the posterior mean of each structural object
  induced by a marginal increase in the weight placed on the indicated likelihood
  component. Larger values indicate that the posterior mean is more responsive to
  that component. For example, the value \(0.50\) for the Control social
  multiplier under Deworming take-up implies that, locally, increasing the weight
  on the take-up likelihood by 10 percent moves its posterior mean by
  approximately \(0.05\) posterior standard deviations. Because the table reports
  absolute values, this example describes the magnitude rather than the direction
  of the response. For objects defined separately by experimental arm, the
  italicized row reports the root mean square of the four arm-specific
  sensitivities. Bold entries indicate the largest sensitivity within each row.
  All distance-dependent objects are evaluated at \(d=1.5\) km.
}
  \end{threeparttable}
\end{table}


\newpage

\begin{table}[ht!]
  \scriptsize
  \caption{Structural Parameters: Posterior Means and 95 Percent Credible Intervals}
  \label{tab:struct-model-params}
  \centering
  \begin{threeparttable}
    \input{tables/all-structural-parameters-table-fit-1-STRUCTURAL_LINEAR_U_SHOCKS_PHAT_MU_REP_FOB}
\floatfoot{\textit{Notes:}
This table reports posterior means and 95 percent credible intervals from the joint Bayesian estimation in Section~\ref{structural}. The model is estimated in Stan using Hamiltonian Monte Carlo with four chains, 400 warm up draws, and 400 saved draws per chain. Sampler diagnostics show no divergent transitions; among the reported parameters, the maximum split \(\hat R\) is 1.011, the minimum bulk effective sample size is 465, and the minimum tail effective sample size is 602. \textit{Take-up parameters.} The parameter \(\lambda\) is the social image weight: it scales the probability that an individual's deworming status is observed into the utility payoff from being seen to deworm. The arm specific private payoff parameters \(\beta_z\), for \(z\in\{\text{Control},\text{Ink},\text{Calendar},\text{Bracelet}\}\), capture non-social utility shifts associated with each assignment arm, including any private item value. The calendar and bracelet private payoff parameters are linked through the relative valuation restriction described in Section~\ref{structural}. The parameter \(\delta\) is the marginal disutility of distance to the treatment site. The parameter \(\sigma_u\) is the standard deviation of the idiosyncratic utility shock \(u_i\). \textit{Observability parameters.} The observability parameters describe the endline belief model for whether an individual's deworming status is known in the community. The parameter \(\beta^O_{\text{Control}}\) is the Control intercept in the observability logit index, while \(\beta^O_z\), for \(z\in\{\text{Ink},\text{Calendar},\text{Bracelet}\}\), are arm specific intercept shifts relative to Control. The parameter \(\gamma^O_{\text{Control}}\) is the Control distance coefficient, while \(\gamma^O_z\), for \(z\in\{\text{Ink},\text{Calendar},\text{Bracelet}\}\), are arm specific distance shifts relative to Control. A negative \(\gamma^O_{\text{Control}}\) means that observability falls with distance in Control communities. Positive values of \(\gamma^O_{\text{Bracelet}}\) or \(\gamma^O_{\text{Ink}}\) indicate that those arms offset the distance related decline in observability. \textit{Relative valuation parameters.} The relative valuation parameters link the bracelet calendar gift choice exercise to private item values. The parameter \(\psi\) is the mean calendar minus bracelet valuation difference in USD, matching the scaled monetary units used in the elicitation likelihood. The parameter \(\rho\) maps that valuation difference into the latent utility scale used in the take-up model. The joint model combines the take-up analysis sample (\(N=9{,}805\)), the relative valuation sample (\(N=467\)), and the endline observability sample (\(N=1{,}141\)).}
  \end{threeparttable}
\end{table}

\begin{table}[ht!]
	\scriptsize
	\caption{Structural Model: Item/Signal Effects on Observability by Distance}
    \label{tab:beliefs-table}
	\centering
    \begin{threeparttable}
    \resizebox{\textwidth}{!}{%
    \input{tables/fit105/fob-beliefs-table}
	}
\floatfoot{\textit{Notes:}
This table reports structural counterparts to the continuous distance reduced form estimates in Table~\ref{tab:rf-fob-tes-cts-dist}. Predictions average over lognormal distance distributions fitted separately for Close and Far communities. Point estimates are posterior means and parentheses contain 95 percent credible intervals. Rows for Bracelet, Calendar, and Ink report effects relative to Control. The Control row reports the model implied Control mean in columns (1)--(3) and the Control Far--Close difference in column (4). The final two rows report posterior differences in expected observability: Any Signal minus No Signal and Bracelet minus Calendar, where Any Signal pools Bracelet and Ink and No Signal pools Control and Calendar. Estimation details, sample sizes, and sampler diagnostics are reported in Appendix Table~\ref{tab:struct-model-params}.}
  \end{threeparttable}
\end{table}


\begin{table}[ht!]
	\scriptsize
	\caption{Structural Model: Item/Signal Effects on Deworming Take-up by Distance}
    \label{tab:struct_overall_effects}
	\centering
    \begin{threeparttable}
    \resizebox{\textwidth}{!}{%
    \input{tables/fit105/struct-overall-te-table}
	}
    \floatfoot{\textit{Notes:}
    This table reports structural counterparts to the continuous distance reduced form estimates in Table~\ref{tab:rf-ates-cts-dist}. Predictions average over lognormal distance distributions fitted separately for Close and Far communities. Point estimates are posterior means and parentheses contain 95 percent credible intervals. Rows for Bracelet, Calendar, and Ink report effects relative to Control. The Control row reports the model implied Control mean in columns (1)--(3) and the Control Far--Close difference in column (4). The final two rows report posterior differences in expected take-up: Any Signal minus No Signal and Bracelet minus Calendar, where Any Signal pools Bracelet and Ink and No Signal pools Control and Calendar. Estimation details, sample sizes, and sampler diagnostics are reported in Appendix Table~\ref{tab:struct-model-params}.}
\end{threeparttable}
\end{table}

\newpage

\begin{table}[ht!]
	\scriptsize
	\caption{Structural Decomposition of Item/Signal Effects on Take-up}
    \label{tab:signaling_private}
	\centering
    \begin{threeparttable}
    \resizebox{\textwidth}{!}{%
    \input{tables/fit105/private-signal-te-table}
	}
    \floatfoot{\textit{Notes:}
    Panel A reports the treatment effect on deworming take-up from social image utility, holding private utility fixed at its Control value. Panel B reports the treatment effect attributable to private utility, holding the social image term fixed at its Control value. Panel B reports pooled private value components only. Point estimates are posterior means and parentheses contain 95 percent credible intervals. Estimation details, sample sizes, and sampler diagnostics are reported in Appendix Table~\ref{tab:struct-model-params}.}
\end{threeparttable}
\end{table}

\begin{table}[ht!]
	\scriptsize
	\caption{Structural Average Treatment Effects: Alternative Observability Measures}
    \label{tab:struct-observability-robustness}
	\centering
    \begin{threeparttable}
    \resizebox{\textwidth}{!}{%
    \input{tables/fit105/struct-observability-robustness-overall-te-table}
	}
    \floatfoot{\textit{Notes:}
    This table reports structural average treatment effects from three variants of the baseline specification that change the endline observability input. Point estimates are posterior means and parentheses contain 95 percent credible intervals. The Correct observability specification uses whether respondents correctly report peers' administrative deworming status as the observability outcome. The Perceived observability specification uses respondents' reports of whether peers know the respondent's own deworming status. The Missing data marginalized specification keeps rows without observed first-order belief responses in the observability margin and averages over their posterior predicted observability probabilities, while observed belief rows contribute to the belief likelihood. Rows for Bracelet, Calendar, and Ink report effects relative to Control. The Control row reports the model implied Control mean in columns (1)--(3) and the Control Far--Close difference in column (4). The final two rows report posterior differences in expected take-up: Any Signal minus No Signal and Bracelet minus Calendar, where Any Signal pools Bracelet and Ink and No Signal pools Control and Calendar.}
\end{threeparttable}
\end{table}

\clearpage
\newpage

\input{tables/density_table}

\vspace{4cm}


\begin{table}[ht!]
	\scriptsize
	\caption{Optimal Number of Points of Treatment under Alternative Social Image Assumptions}
    \label{tab:optim-post}
	\centering
    \begin{threeparttable}
    \resizebox{\textwidth}{!}{%
    \input{tables/fit105/optim-summ-table}
	}
    \floatfoot{\textit{Notes:}
    Point estimates are posterior medians across 1{,}600 draws and parentheses contain 95 percent credible intervals. Coverage, mean take-up, and mean one-way distance use census adult-population weights for the 144 communities. Within each draw, all scenarios use the same target: the expected number of adults treated under the experimental allocation and Control observability. In Panel A, site counts and mean distance are fixed, so their intervals are omitted. Panel B solves equation~\eqref{eq:oa-problem} using all 1{,}451 candidate sites and a 3.5 km distance cap. The Control and Bracelet rows allow social-image returns to vary with distance; the fixed-return rows hold the corresponding returns at their 0.5 km values. The no-social-image row sets \(\mu=0\). \textsuperscript{$\dagger$}Its target is infeasible in every draw, so it reports the best feasible allocation. Panel C reports paired Control--Bracelet site savings and the additional savings from endogenous returns relative to fixed returns, computed within each draw. These need not equal differences between marginal posterior medians.}
\end{threeparttable}
\end{table}


\clearpage

\FloatBarrier



\subsection{Baseline Sample Balance, Demographics, Deworming Knowledge, and Social Image Concerns}
\label{app:baseline-sample}

Table~\ref{tab:baseline-bal} reports baseline characteristics and balance tests for the 1{,}995 respondents in the baseline sample. Panel A reports demographic characteristics, Panel B reports deworming knowledge and prior deworming experience, and Panel C reports baseline social image measures. For each variable, the table reports no-item Control means and pairwise differences by item/signal arm, separately for Close and Far communities, along with joint tests of equality across arms.

  Panel B shows some imbalance in baseline deworming knowledge and prior deworming
  experience. In particular, the joint tests are significant at the ten percent
  level for three measures: knowledge that medication treats worms, deworming in the
  past year, and knowledge that biannual treatment is recommended. Because baseline
  survey records cannot be linked at the individual level to the monitored
  take-up data, we construct community-level baseline
  shares for these three measures and include them as controls in the main reduced-form specifications. The results are reported in Tables~\ref{tab:fob-baseline-imbalance-robustness} and \ref{tab:prior-deworming-robustness}. The key patterns
  are stable: for reported knowledge of peers' deworming status, the Bracelet and
  Ink Far--Close interactions are 18.2 and 11.3 percentage points, while the
  Calendar interaction is 2.8 percentage points; for administrative take-up, the
  corresponding interactions are 12.1, 10.7, and 2.9 percentage points. Thus, the
  central observability and take-up patterns are not explained by these baseline
  survey imbalances.

  Overall, the baseline balance exercise shows no systematic pattern of imbalance across item/signal arms. Statistically significant differences are limited to the Panel B deworming knowledge and prior experience measures discussed above, and the main results are robust to controlling for community-level baseline shares of these variables. The social image measures in Panel C are balanced across arms. The table also documents several features of the setting used in Section~\ref{setting}. Deworming knowledge is widespread: 93.4 percent of respondents know that children can get worms, 69.9 percent know that adults can get worms, and 78.2 percent know that medication treats worms. Prior experience with deworming is common, with 68.4 percent reporting ever having dewormed and 37.4 percent reporting deworming in the past year. Baseline social image concerns are also strong: 94.7 percent of respondents say they would praise someone who deworms, and 69.9 percent say they would judge someone who does not.


% Overall, the baseline sample is broadly balanced across item/signal arms. Statistically significant differences are sparse relative to the number of reported comparisons and do not reveal a systematic pattern across treatment arms. In particular, the social image measures in Panel C are balanced across arms. The table also documents several features of the setting used in Section~\ref{setting}. Deworming knowledge is widespread: 93.4 percent of respondents know that children can get worms, 69.9 percent know that adults can get worms, and 78.2 percent know that medication treats worms. Prior experience with deworming is common, with 68.4 percent reporting ever having dewormed and 37.4 percent reporting deworming in the past year. Baseline social image concerns are also strong: 94.7 percent of respondents say they would praise someone who deworms, and 69.9 percent say they would judge someone who does not.

\begin{table}[ht!]
	\scriptsize
	\caption{Baseline Sample Balance: Demographics, Deworming Knowledge, and Social Image Concerns}
    \label{tab:baseline-bal}
	\centering
    \begin{threeparttable}
    \resizebox{\textwidth}{!}{%
    \input{tables/baseline-demo-deworming-knowledge-balance-table}
	}
       \floatfoot{\begin{minipage}{\linewidth}\textit{Notes:}
       The table reports baseline characteristics for the 1{,}995 respondents in the baseline survey. Panel A reports demographic characteristics, Panel B reports deworming knowledge and prior deworming experience, and Panel C reports social image measures. The Sample Mean column reports the mean in the baseline sample. Columns labeled Con report the no-item Control mean, with standard errors in parentheses. Columns labeled Ink - Con, Cal - Con, Bra - Con, and Bra - Cal report differences relative to the indicated comparison group, with \(p\)-values in square brackets below. Joint Test reports the \(F\)-test \(p\)-value for equality of means across item/signal arms within each distance condition and, in the final column, across both distance conditions combined. Standard errors are clustered at the community level, and county fixed effects are included to reflect stratification.

       Variable definitions:

    \begin{itemize}\setlength\itemsep{0.2em}
  \item ``Know adults get worms'' equals one if the respondent lists adults among those who can be infected.
  \item ``Know children get worms'' equals one if the respondent lists children among those who can be infected.
  \item ``Believe deworming is for the sick'' equals one if the respondent states that only sick people should deworm.
  \item ``Know medication treats worms'' equals one if the respondent states that medicine cures worm infections.
  \item ``Dewormed in the past'' equals one if the respondent reports ever deworming.
  \item ``Dewormed in the past year'' equals one if the respondent reports deworming in the last 12 months.
  \item ``Know bi-yearly treatment recommended'' equals one if the respondent correctly identifies the recommended treatment frequency.
  \item ``Know worms impose externalities'' equals one if the respondent answers ``Yes'' to ``Can a person sick with worms spread worms to others?'' or to both of the following: ``If you have worms, does that affect your neighbors' or relatives' health?'' and ``If your neighbors or relatives have worms, does that affect your health?''
  \item ``Would you praise: Deworming?'' equals one if the respondent says they would praise someone who came for free deworming treatment at a central location.
  \item ``Would you judge: Not deworming?'' equals one if the respondent says they would look down on someone who did not come for free deworming treatment at a central location.
  \item The child immunization praise and judgment variables are defined analogously for the child immunization vignette.
\end{itemize}
\end{minipage}}
\end{threeparttable}
\end{table}

\newpage

\newcommand{\samplesubhead}[1]{%
  \par\addvspace{0.45\baselineskip}%
  \noindent\emph{#1}\par\nopagebreak\noindent
}

\subsection{Sample Construction}
\label{app:sample-construction}

Several independently drawn rosters enter the final analysis samples because the field experiment combined a large community deworming campaign with census based sampling, SMS recruitment, administrative take-up monitoring, and endline surveying. Figure~\ref{fig:sample-consort} displays how these rosters map onto the main analysis samples.

\samplesubhead{Household census.}
Enumerators completed a household census in the 144 implemented clusters and recorded 39{,}301 adults aged 18 or older. The census serves as the master sampling frame for the baseline survey, SMS recruitment, no-SMS rosters, endline survey, and administrative take-up monitoring. Appendix~\ref{sec:site-selection} describes the cluster selection process and the field constraints that led to implementation in 144 clusters.

\samplesubhead{Baseline survey.}
We constructed a baseline roster by randomly sampling up to 17 households per cluster and allowing up to two adults per household, in order to have backup respondents when sampled adults could not be reached. This generated a baseline roster of 2{,}414 households and 4{,}185 adults across the 144 implemented clusters.\footnote{The number of sampled households is slightly below \(17 \times 144 = 2{,}448\) because the realized household draw was smaller in a small number of clusters.} Enumerators then attempted to reach the prespecified target of 2{,}160 adults for baseline surveys. Of the 2{,}062 adults successfully reached, 38 were unable to participate, for example due to illness or language barriers, and 29 refused consent. This yielded 1{,}995 completed baseline surveys, which constitute the Baseline sample.

The full 4{,}185 person baseline roster also enters the administrative monitoring lists used to record deworming take-up at points of treatment during the campaign. Thus, the baseline roster contributes to the broader monitored take-up sample described below.

\samplesubhead{SMS treatment and no-SMS rosters.}
Using the census roster, we randomly sampled phone-owning adults for SMS recruitment. In Control clusters, we sampled 17 phone-owning adults for Reminder Only and 17 for Social Information. In clusters assigned to Bracelet, Ink, or Calendar, we sampled 29 phone-owning adults for Social Information. Where possible, we also drew within-household backup phone owners to use when the initially sampled individual could not be reached. Only one adult per household could ultimately be enrolled in an SMS treatment. Enumerators recruited sampled individuals in person before the campaign. They reached 3{,}040 adults, of whom 18 did not consent. The final enrolled SMS treatment sample consists of 3{,}022 adults: 387 Reminder Only recipients and 2{,}635 Social Information recipients. Only adults who were reached, consented, and enrolled in the SMS service were added to the take-up monitoring list.

To form the phone owner no-SMS comparison group used in the SMS mechanism analysis, we randomly sampled phone owners from households not selected for SMS treatment. The primary phone owner no-SMS comparison group contains 3{,}813 adults. These individuals received no SMS messages and no recruitment visit. The SMS mechanism analyses compare the enrolled SMS treatment sample to this phone owner no-SMS comparison group. Appendix Tables~\ref{tab:sms-balance} and~\ref{tab:sms-enrollment} report balance and enrollment checks for these SMS samples, and Section~\ref{sec:sms-alternatives} reports the SMS mechanism results.

For administrative take-up monitoring and endline sampling, we also constructed broader no-SMS rosters. For phone owners, an additional within-household phone owner was sampled where possible, yielding a broader phone owner no-SMS roster of 5{,}803 adults. For non-phone owners, we sampled 8 non-phone owners per cluster, yielding 1{,}138 adults, and added a backup draw of 576 adults. Where possible, an additional non-phone owner was sampled within the same household. After deduplication, the intended non-phone owner no-SMS roster contains 2{,}341 adults. Combining the broader phone-owner roster and the intended non-phone owner roster yields 8{,}144 no-SMS adults. Together with the 3{,}022 enrolled SMS recipients, these adults form the 11{,}166 person SMS/endline roster used for endline sampling.

\samplesubhead{Administrative take-up data.}
Deworming take-up was recorded administratively at points of treatment for individuals on the monitoring lists. These lists were assembled from three independently drawn subsamples of the census roster: the enrolled SMS treatment sample of 3{,}022 adults, the no-SMS monitoring rosters, and the baseline roster of 4{,}185 adults.

A SurveyCTO list-generation error caused 989 wave-1 non-phone owner no-SMS controls to be omitted from the administrative monitoring lists. This error affected administrative take-up monitoring, not the SMS mechanism comparison or the endline sampling frame. As a result, the monitored non-phone owner no-SMS sample contains 1{,}352 adults rather than the intended 2{,}341. Combining the 5{,}803 phone owner no-SMS adults with the 1{,}352 monitored non-phone-owner no-SMS adults yields 7{,}155 no-SMS adults in the administrative monitoring lists.

Because the enrolled SMS treatment sample, no-SMS monitoring rosters, and baseline roster were drawn independently from the census, some individuals appear on more than one list. After deduplication, the combined administrative monitoring universe comprises 12{,}827 unique adults. Our main reduced-form analyses of the distance and item/signal treatments exclude the 3{,}022 adults enrolled in SMS treatments. This yields the main take-up analysis sample of 9{,}805 adults.

\samplesubhead{Endline survey.}
The endline sampling frame was the 11{,}166 person SMS/endline roster, consisting of enrolled SMS recipients and no-SMS adults. From this frame, we randomly selected a main endline roster of 3{,}586 adults, consisting of 880 SMS treated adults and 2{,}706 no-SMS adults, and added 634 backup respondents, consisting of 350 SMS treated backups and 284 no-SMS backups. The resulting endline field roster contained 4{,}220 adults. Enumerators successfully reached 3{,}830 of them. Of those reached, 78 were unable to participate and 23 refused consent, yielding 3{,}729 completed endline surveys.

The prespecified target was 3{,}600 completed endline surveys, equal to 25 per cluster across 144 clusters. Completed surveys exceeded this target because backup individuals were systematically surveyed during implementation. Excluding adults enrolled in the SMS treatment arms, which are analyzed separately, yields the Endline sample of 2{,}659 respondents used for the main distance and item/signal analyses.

The endline survey generates several additional measures used throughout the analysis. First, 2{,}039 endline respondents were assigned to an observability module eliciting knowledge of peers' deworming take-up. Of these, 1{,}784 completed the module and 255 have missing observability module data due to a programming error, discussed in Section~\ref{sec:sample-balance}. After excluding adults enrolled in SMS treatments, the Endline observability sample contains 1{,}141 respondents.

Second, all endline respondents chose between a bracelet and a calendar as a thank-you gift after completing the survey. The main gift-choice analysis focuses on adults who did not deworm, our closest proxy for marginal participants, because they did not receive an item through the campaign. The full non-dewormer gift-choice sample includes 1{,}350 respondents. Excluding adults enrolled in the SMS treatment arms yields the no-SMS gift-choice analysis sample of 1{,}174 respondents used in Table~\ref{tab:pref-cal-gift}.

Third, in no-item Control clusters, we conducted a relative valuation exercise for the bracelet and calendar. Respondents first chose between the two items as a gift. We then elicited their willingness to accept compensation to give up their preferred item for the other. This exercise yields a measure of the monetary value gap between the bracelet and calendar. We attempted to reach 600 individuals across 34 Control clusters, successfully reached 474, and obtained completed relative valuation data from 467 respondents. We refer to these respondents as the Relative valuation sample.

\begin{landscape}
    \begin{figure}
      \centering
      \adjustbox{max width=\linewidth}{\input{tables/manual/sample-consort}}
      \caption{Sample Construction and Analysis Samples}
      \label{fig:sample-consort}
\floatfoot{\textit{Notes:} Green boxes denote analysis samples used in the main text. The administrative take-up data combine individuals drawn from the baseline roster and the enrolled SMS and no-SMS monitoring rosters, deduplicated across rosters. Because these rosters were drawn independently from the census, individuals may appear in more than one roster before deduplication. The main distance and item/signal analyses exclude adults enrolled in the SMS treatments, which are analyzed separately.}
    \end{figure}
\end{landscape}

\newpage

\subsection{PoT Selection and Distance Randomization}\label{sec:site-selection}

This appendix provides implementation details on the geographically constrained selection of PoT-centered clusters, the targeted community selection, and the Close/Far distance classification. The notation follows the cluster selection algorithm in Section~4.4 of the pre-analysis plan.

\paragraph{Definitions.}
A point of treatment (PoT) is the finalized location at which CHVs delivered deworming during the campaign. Because geocoded community boundaries were not available ex ante, we used geo-coded primary schools as spatial reference points for constructing candidate PoT catchments and guiding field listing. We refer to these school-based locations as candidate PoT centers. Candidate PoT centers are used for sampling and assignment and need not coincide with the finalized PoT. During field mapping, feasible PoT locations were identified near the assigned reference points, and final distance classifications were verified using finalized community-centroid-to-PoT distances. A candidate cluster is defined by a candidate PoT center together with its surrounding 2.5~km catchment radius. This 2.5~km catchment delimits the set of communities eligible to be selected as the targeted community for that PoT.

\paragraph{Sampling frame.}
We begin from a sampling frame of 1{,}451 geocoded primary schools across the three study counties, which serve as candidate PoT centers.\footnote{Coordinates are from the Kenya Open Data Portal: \url{http://www.opendata.go.ke/}.} Each school defines a candidate PoT-centered cluster with a 2.5~km catchment radius. We drew 150 target clusters and eight fallback clusters, for a total of 158 candidate clusters. Field feasibility constraints limited implementation to 144 clusters. Eight candidate PoTs were dropped during mapping because the school's location could not be verified (two in Busia, two in Siaya, and four in Kakamega). An additional six clusters from the original target list were dropped because of implementation constraints, including lack of available CHVs, opposition from local authorities, or community resistance (two in Siaya, two in Busia, and two in Kakamega). Attrition was higher among Far clusters (8 vs.\ 6), but this difference is consistent with chance under random assignment (binomial test, \(p=0.79\)). Appendix Table~\ref{tab:stacked-sample-balance-table} shows that covariates remain balanced in the final sample.


\subsubsection*{PoT selection under spatial overlap constraints}

We select clusters using a sequential acceptance--rejection algorithm. Let \(C\) denote the set of all candidate clusters (\(|C|=1{,}451\)) and let \(R\) denote the set of accepted clusters. The algorithm iterates through candidates drawn at random, adding a candidate cluster to \(R\) only if doing so satisfies a minimum non-overlapping catchment-size requirement for all clusters in \(R\). If the algorithm cannot reach the target number of clusters under the constraints, the procedure restarts from the full candidate set.

\paragraph{Non-overlapping catchment size.}
For a candidate cluster \(i\) and an accepted set \(R\), define the non-overlapping catchment of \(i\) as the portion of \(i\)'s 2.5~km catchment circle that does not lie within the larger buffer areas used for overlap checks around already selected PoTs. We operationalize the size of this non-overlapping catchment by counting available primary school neighborhoods. Specifically, a school neighborhood counts as available if at least 60\% of the 1~km radius circle around the school lies within the non-overlapping portion of the 2.5~km catchment. We require each selected cluster to retain at least one such available school neighborhood, i.e., \(\text{ThresholdClusterSize}=1\). Intuitively, this ensures that overlap restrictions do not leave a cluster with an effectively empty area from which to list eligible communities.

\paragraph{Buffers used for overlap checks.}
Although the 2.5~km radius defines the eligibility region for selecting study communities, overlap restrictions are evaluated using larger buffers to create separation between nearby PoTs. For the purpose of the site selection algorithm, candidate PoTs are assigned a buffer radius of either 3~km or 4~km. Figure~\ref{fig:cluster-overlap-schematic} illustrates the geometry: the shaded region represents the portion of a selected PoT's 2.5~km catchment that is excluded because it falls within a neighboring candidate PoT's buffer. Study communities are selected from the remaining, non-overlapping portion of the 2.5~km catchment. As a result, in the candidate-center geometry used by the assignment algorithm, a selected study community is separated from alternative candidate PoT catchments by the imposed overlap restrictions.


\subsubsection*{Randomized anchor selection and distance targeting}

Within each selected cluster, the algorithm identifies eligible primary schools that can serve as anchors for field listing. Each potential anchor is classified according to its distance from the candidate PoT center. Anchors within 1.25~km of the candidate center lie in the inner band, while anchors between 1.25 and 2.5~km lie in the outer band. The 1.25~km cutoff was chosen ex ante as the midpoint of the 2.5~km catchment radius.

One eligible anchor is then selected at random. The randomization is stratified by county, with selection probabilities designed to produce equal numbers of inner- and outer-band anchor selections. The selected anchor determines the area in which enumerators search for and list eligible communities. Thus, the anchor randomization shifts the distribution of distances at which targeted communities are likely to be located.

Importantly, the selected anchor's band is a randomized spatial sampling target, not the Close/Far classification of the community ultimately included in the study. Because community centroids were unavailable at the time of the anchor draw, the design could target field listing toward the inner or outer portion of the catchment but could not directly assign the subsequently selected community to a realized distance category. The dotted circle in Figure~\ref{fig:cluster-overlap-schematic} depicts the 1.25~km boundary used both to define the two anchor bands ex ante and to classify realized community distance ex post.

\subsubsection*{Targeted community selection and realized distance classification}

Enumerators listed eligible communities in the vicinity of the randomly selected anchor, and one community was randomly selected from this listing as the cluster's targeted community. As part of field mapping, the final feasible PoT location was identified near the candidate PoT center and the centroid of the selected community was recorded. We then calculated the distance between the community centroid and the finalized PoT. Communities with realized distances of at most 1.25~km were classified as Close, while those with realized distances greater than 1.25~km and at most 2.5~km were classified as Far.

The realized classification need not coincide with the distance band of the selected anchor. The anchor school is only a proxy for the location of the subsequently selected community, and the finalized PoT may differ from the candidate PoT center used by the spatial algorithm. Consequently, an inner-band anchor can yield a community whose realized distance is Far, and an outer-band anchor can yield a community whose realized distance is Close. These cases represent differences between the randomized distance target and realized community distance, rather than corrections to an initially assigned community treatment.

All community mapping and realized distance classification occurred before implementation and before any outcome data were observed. In the realized implementation data, each selected community's assigned finalized PoT is the closest finalized deworming site.

\subsubsection*{Randomized anchor targets and realized distance summary}

Among the 158 selected candidate clusters, the anchor-selection procedure produced 80 inner-band and 78 outer-band anchor targets. The one-cluster deviation from an even 79--79 split reflects an off-by-one error in the original randomization algorithm. Fourteen clusters were dropped before any outcome data were observed: six with inner-band anchor targets and eight with outer-band anchor targets.

Among the 144 implemented clusters, the realized distance category differed from the randomized anchor target in 26 cases. Sixteen communities listed using outer-band anchors had realized distances within the Close band, while ten communities listed using inner-band anchors had realized distances within the Far band. The final implemented sample therefore contains 80 communities classified as Close and 64 classified as Far. Under a null of equal allocation, the two-sided binomial \(p\)-value for an imbalance at least this large is \(0.211\), so the null of equal allocation cannot be rejected.


\input{figures/tikz/anchor-diagram}

\newpage

\subsubsection*{Potential distance balance under constrained PoT selection}

A potential concern is that the acceptance-rejection algorithm may have introduced systematic differences in the set of feasible distances for communities classified as Close versus Far using realized distance, complicating the interpretation of realized continuous distance. Following \citet{borusyak2023nonrandom}, we assess this by simulating the realized selection and assignment mechanism. Holding the set of surveyed communities fixed, we re-run the full PoT selection, anchor selection, targeted-community pairing, and feasibility-screening algorithm 500 times. In each simulated draw, each realized surveyed community receives the centroid-to-PoT distance it would have had under the feasible simulated PoT-community pairing in which it is selected. We then compare the resulting potential distance distributions for communities in the realized Close and Far distance categories. Figure~\ref{fig:bh-dens-plot} shows that these distributions are statistically indistinguishable, providing no evidence that the constrained algorithm generated systematic differences in potential distance profiles across realized distance categories.


\begin{figure}[htbp]
  \centering
  \includegraphics[width=4.8in]{misc-figures/counterfactual-assignment-density-close-far-dens.pdf}
  \caption{Potential Distance Distributions under Constrained Randomization}
  \label{fig:bh-dens-plot}
\floatfoot{\textit{Notes:} The figure plots potential community-centroid-to-PoT distance distributions implied by 500 re-runs of the full PoT selection and assignment algorithm \citep{borusyak2023nonrandom}, holding fixed the set of surveyed communities. Red (blue) curves correspond to communities classified as Close (Far) in the realized experiment. The plotted unit is the realized surveyed community by simulation draw, weighted equally, and the distances are counterfactual potential distances under simulated feasible PoT-community pairings rather than the observed community-centroid-to-PoT distances. A Kolmogorov--Smirnov test comparing these community-draw potential-distance distributions fails to reject equality (\(p=0.501\)), indicating no evidence that Close and Far communities differ in their potential distance profiles under the constrained randomization.}
\end{figure}


\clearpage




\FloatBarrier


\subsubsection{Continuous Distance Randomization Inference}
\label{sec:distance-random}

As an additional check that our continuous distance measure is orthogonal to pretreatment characteristics, we regress each baseline covariate on the distance from the community centroid to its assigned point of treatment (PoT), controlling for county fixed effects and clustering standard errors at the community level. We then use the constrained assignment simulations described above to construct placebo continuous-distance draws. For each placebo draw, each surveyed community receives a centroid-to-PoT distance sampled from the feasible PoT-community pairings generated for that community by re-running the PoT selection, Close/Far assignment, targeted-community pairing, and feasibility-screening algorithm. We then recompute the associated test statistic. Figure~\ref{fig:ri-baseline} plots the resulting placebo distribution for each covariate under the null of no association. The vertical line marks the realized value, and simulation-based \(p\)-values based on the 500 draws are reported in the upper right-hand corner of each panel. For all covariates shown, the \(p\)-value exceeds 0.10, indicating that realized centroid-to-PoT distance is not systematically related to observable characteristics.

\begin{figure}[htbp]
  \centering
  \includegraphics[width=5.2in]{figures/dist-ri-plot.pdf}
  \caption{Randomization Inference: Baseline Covariates and Centroid-to-PoT Distance}
  \label{fig:ri-baseline}
   \floatfoot{\textit{Notes:} Each histogram shows the placebo distribution of the OLS \(t\)-statistic from 500 counterfactual continuous-distance draws based on simulated feasible PoT-community pairings under the constrained assignment algorithm. Regressions include county fixed effects. Standard errors are clustered by community. The vertical line is the realized statistic from the actual experiment, and each panel reports the corresponding simulation-based \(p\)-value.}
\end{figure}

\FloatBarrier
\newpage


%---------------------------------------------
% PAP mapping table (multi-page, appendix)
%---------------------------------------------

\subsection{Pre-Analysis Plan: Deviations and Clarifications}
\label{sec:pap}

This experiment was preregistered with the AEA RCT Registry as AEARCTR-0001643. This appendix documents analyses added after the pre-analysis plan (PAP), implementation and sampling clarifications, and data-collection issues that affect the construction of analysis samples. Table~\ref{tab:pap_mapping} maps preregistered items to the corresponding analyses in the paper.

\paragraph{Post-PAP additions.}
We added three analyses or measurements after the PAP. First, we develop and estimate a structural model of social image, building on \citet{Benabou_laws_2011}, to quantify the social-image multiplier. Second, we use the estimated model in a site allocation exercise. Third, in Control communities, we conducted an endline relative valuation elicitation comparing the bracelet and calendar: respondents first chose between the two items as a gift, and we then elicited their willingness to accept compensation to give up their preferred item for the other. The relative valuation elicitation complements the preregistered endline gift choice module. These additions are presented alongside the preregistered reduced form analyses and do not replace them.

\paragraph{Presentation of auxiliary analyses.}
The PAP specifies SMS comparisons involving Reminder Only and Social Information messages. We report these analyses as auxiliary tests of reminder and social information channels rather than as part of the main reduced form evidence on distance and item/signal arms. The main reduced-form analysis compares Close and Far communities using the classification described in Section~\ref{sec:distance-treatment}. Continuous distance enters the reduced-form robustness checks and structural model.

\paragraph{Cluster implementation and distance assignment.}
The initial field plan targeted 150 clusters, with eight additional fallback clusters. Field feasibility constraints limited final implementation to 144 clusters, as reflected in the registered PAP. Eight candidate PoTs were dropped during mapping because the school's location could not be verified, and six additional clusters were dropped because of implementation constraints, including lack of available CHVs, opposition from local authorities, or community resistance. The selected anchor determines the area in which enumerators search for and list eligible communities. We then calculated the distance between the community centroid and the finalized PoT and applied the prespecified 1.25 km cutoff to classify communities as Close or Far. Appendix~\ref{sec:site-selection} provides details on cluster selection and dropped clusters.

\paragraph{Sample construction.}
Relative to the PAP, three sample construction changes occurred. First, the baseline roster was expanded to allow backup respondents, while enumerators still targeted the prespecified number of baseline interviews. Second, more phone-owning adults were sampled for SMS recruitment than specified in the PAP, in order to preserve recruitment targets in the presence of non-response. Only reached and consenting adults received SMS messages and entered the enrolled SMS treatment sample. Third, backup endline respondents were systematically surveyed, yielding 3{,}729 completed endline surveys relative to the prespecified target of 3{,}600. Appendix~\ref{app:sample-construction} provides the full construction of the baseline, SMS, endline, and administrative monitoring samples.

\paragraph{Data-collection issues affecting analysis samples.}
Two programming or list generation errors affect the construction of the analysis samples. First, 255 respondents assigned to the observability module did not receive it because of a programming error. Section~\ref{sec:sample-balance} discusses this issue, Table~\ref{tab:endline-obs-attrit} reports differential missingness checks, and Table~\ref{tab:rf-fob-lee} reports Lee bounds for the observability outcomes. Second, a SurveyCTO list generation error caused 989 wave 1 non-phone-owning adults in the no-SMS comparison sample to be omitted from the PoT monitoring lists. The no-SMS comparison sample spans all item/signal arms and is distinct from the no-item Control arm. These adults are therefore absent from the administrative take-up data. Section~\ref{sec:sample-balance} discusses this issue, Table~\ref{tab:takeup-attrit} tests whether omission differs across item/signal arms, and Appendix~\ref{app:sample-construction} describes the construction of the administrative monitoring sample.

\paragraph{Other preregistered diagnostics.}
The PAP also specified several diagnostic and auxiliary analyses, including SMS airtime reward checks, temporal dynamics of take-up, sample selection adjustments related to phone ownership, inter-cluster spillover checks, and multiple hypothesis correction. Table~\ref{tab:pap_mapping} lists these items and indicates whether they are reported in the paper, reported in the appendix, or not used for the paper's main claims.


\begin{landscape}
\input{tables/manual/pap-mapping}
\end{landscape}
\newpage

\FloatBarrier



\subsection{Perceived Meaning and Negative Perceptions of Items/Signals}
\label{app:bad_incentive}

This appendix provides descriptive evidence on how respondents interpreted the assigned items/signals and whether they expressed negative associations with them. The sample is restricted to endline respondents in the no-SMS group and to the item/signal arms Calendar, Ink, and Bracelet.

\paragraph{Perceived meaning and stigma.}
Respondents were asked whether they had seen the assigned item or signal in the community and, conditional on reporting that they had seen it, what it meant if someone had that item or signal. Using the same deworming-linked indicator as in Table~\ref{tab:incentive-check}, among respondents who report seeing the item or signal and provide a non-missing meaning response (\(N=1{,}669\)), the overwhelming majority interpret it as indicating deworming participation or treatment: 96.6 percent in the Ink arm, 94.6 percent in the Bracelet arm, and 89.3 percent in the Calendar arm mention deworming, medication, treatment, tablets, drugs, or worms. Explicitly stigmatizing interpretations are essentially absent in these meaning responses: none of the 1{,}669 responses contain terms suggesting that the person is dirty, unclean, infected, or sick. An illustrative deworming linked response is: ``It shows the person went for deworming treatment.''

\paragraph{Bracelet visibility and wearing by distance.}
Table~\ref{tab:bracelet-distance-display} examines whether bracelet visibility, retention, or wearing differs by distance. Bracelet visibility is high in both Close and Far communities: 95.2 percent of respondents in Close communities and 93.8 percent in Far communities report seeing the bracelet in the community. Among respondents who dewormed and report receiving a bracelet, 76.2 percent in Close communities and 86.1 percent in Far communities still had it at endline, while 48.2 percent and 42.9 percent, respectively, were wearing it. Thus, bracelets were highly visible in both distance categories, and there is no evidence that Far respondents were more likely to actively display the bracelet by wearing it.


\begin{table}[ht!]
    \scriptsize
    \caption{Bracelet Visibility, Retention, and Wearing by Distance}
    \label{tab:bracelet-distance-display}
    \centering
    \begin{threeparttable}
    \input{rf-tables/main-specs/bracelet-distance-display-tbl}
    \floatfoot{\textit{Notes:} The table reports endline measures of bracelet visibility, possession, and wearing among respondents assigned to the Bracelet arm. Outcomes are reported separately for Close and Far communities, along with Far--Close differences. ``Seen in community'' equals one if the respondent reports having seen the bracelet in the community. ``Still has bracelet'' and ``Wearing bracelet'' are defined among respondents who dewormed and report receiving the bracelet. Standard errors, clustered at the community level, are reported in parentheses.}
    \end{threeparttable}
\end{table}

\paragraph{Reported drawbacks.}
To assess whether the items/signals generated stigma or discomfort, endline respondents were asked whether there was anything they did not like about the assigned item or signal. Respondents answering ``yes'' were asked an open-ended follow-up explaining what they did not like. These drawback questions were collected only for Wave 2 respondents, which is why the number asked is smaller than the full endline sample. Table~\ref{tab:incentive-check-backdraw} reports the share of respondents reporting any drawback by item/signal arm. Complaints are uncommon overall: the Ink mean is about 3.3 percent, and the Bracelet complaint rate is statistically indistinguishable from Ink. Calendar recipients are more likely to report some drawback, about 5.5 percentage points higher than Ink, but the open-ended explanations indicate that these concerns primarily reflect design or informational preferences rather than stigma.

\paragraph{Content of open-ended drawbacks.}
Table~\ref{tab:bad_breakdown_combined} codes the open-ended drawback responses into salient themes. For the calendar, the modal complaint among the small set of complainers is that the calendar should more clearly state who provided it and that it is linked to deworming: 20 of 31 complainers ask for an NGO name, logo, or explicit deworming-campaign message. This pattern is inconsistent with respondents viewing deworming as something to hide; instead, it suggests that some respondents preferred the calendar to communicate the campaign purpose more clearly.

For ink, complaints primarily concern the mark being dirty, sticky, or slow to fade. For bracelets, the most common complaints concern religious objections, fit or comfort, and appearance. Explicitly stigma related comments are rare in this module: one bracelet response associated the bracelet with prostitution, and one calendar response referenced a family planning rumor.

\begin{table}[ht!]
    \scriptsize
    \caption{Reported Drawbacks of Assigned Items/Signals at Endline}
    \label{tab:incentive-check-backdraw}
    \centering
    \begin{threeparttable}
    \input{rf-tables/main-specs/incentive-drawback-tbl}
    \floatfoot{\textit{Notes:}
    The sample is restricted to endline respondents in the no-SMS group, in the Calendar, Ink, or Bracelet arms, who dewormed and were asked the Wave 2 drawback questions. The dependent variable equals one if the respondent reports disliking something about the assigned item or signal and zero otherwise. The sample includes only observations with a non-missing response to the relevant drawback question. Ink is the omitted group. Standard errors clustered at the community level are reported in parentheses. \(^{*}p<0.10\), \(^{**}p<0.05\), \(^{***}p<0.01\).
    }
    \end{threeparttable}
\end{table}

\begin{table}[!htbp]
    \scriptsize
    \centering
    \caption{Open-Ended Drawbacks of Assigned Items/Signals}
    \label{tab:bad_breakdown_combined}
    \begin{threeparttable}
    \resizebox{\textwidth}{!}{%
    \begin{tabular}{lccc}
    \toprule
    Drawback type & Count & Share of complainers & Share of asked respondents \\
    \midrule
    \multicolumn{4}{l}{\textbf{Panel A: Calendar} (31 complainers; 354 asked)} \\
    Low quality/design & 8 & 25.8\% & 2.3\% \\
    Missing NGO/campaign branding or explicit deworming messaging & 20 & 64.5\% & 5.6\% \\
    Other, including family-planning rumor or preference for other reward & 3 & 9.7\% & 0.8\% \\
    \midrule
    \multicolumn{4}{l}{\textbf{Panel B: Ink} (11 complainers; 337 asked)} \\
    Dirty/sticky/spoils hands & 7 & 63.6\% & 2.1\% \\
    Slow to fade / hard to remove & 4 & 36.4\% & 1.2\% \\
    Design/alternative suggested & 2 & 18.2\% & 0.6\% \\
    Fear/concern & 1 & 9.1\% & 0.3\% \\
    \midrule
    \multicolumn{4}{l}{\textbf{Panel C: Bracelet} (14 complainers; 377 asked)} \\
    Religious/church disapproval & 5 & 35.7\% & 1.3\% \\
    Fit/size/comfort & 4 & 28.6\% & 1.1\% \\
    Appearance/color, including gender or age suitability & 5 & 35.7\% & 1.3\% \\
    Misinterpreted as voting & 1 & 7.1\% & 0.3\% \\
    Stigma association: prostitution & 1 & 7.1\% & 0.3\% \\
    \bottomrule
    \end{tabular}
    }
    \floatfoot{\textit{Notes:}
    Drawback themes are coded from open-ended responses among respondents who answered ``yes'' to the item-specific drawback question. The ``asked'' denominator is the number of respondents for whom the corresponding drawback question was observed; these questions were collected only for Wave 2 respondents. \textbf{Calendar} categories are mutually exclusive: \emph{Low quality/design} captures complaints about paper quality, size, and aesthetics; \emph{Missing NGO/campaign branding or explicit deworming messaging} captures requests for an organization name/logo, explicit reference to deworming or worms, and/or campaign logistics; and \emph{Other} includes a family-planning rumor and preferences for other rewards. \textbf{Ink} and \textbf{Bracelet} themes are multi-coded, so a response may contribute to multiple rows.
    }
    \end{threeparttable}
\end{table}


\clearpage


\subsection{Elicitation Format Checks for Reported Observability}
\label{app:demand_checks}

A concern with survey elicited beliefs about peers' deworming participation is that respondents may feel pressure to provide a definite answer rather than responding ``Don't know'' or to overstate peers' participation. This appendix uses the PAP specified randomized elicitation format in the endline social knowledge module, the module from which the main reported observability measure is constructed, to assess these concerns.

\paragraph{Elicitation formats.}
Endline respondents were randomized into one of two formats. In the named individual format, respondents were shown ten named adults from their community one at a time. For each name, they first indicated whether they recognized the person and, if so, whether they thought the person came for deworming: Yes, No, or Don't know. This is the format used to construct the main observability measure in Table~\ref{tab:beliefs-visibility}.

In the paired name format, respondents were shown ten pairs of names. For each pair, they first reported how many of the two people they recognized, without indicating which name or names they recognized. If exactly one person was recognized, respondents were asked whether the person they knew came for deworming, again without revealing which person they knew. If both people were recognized, respondents were asked whether both came for treatment. We exclude these both recognized prompts from the participation belief comparisons because the response cannot be mapped cleanly to an individual level belief: a ``No'' response does not reveal which person, or whether both people, the respondent believes did not come. We therefore use paired name prompts with exactly one recognized person for the Yes/No/Don't know and accuracy comparisons, where the response refers to a single known individual while preserving anonymity.

\paragraph{Checks.}
We conduct three checks. First, we compare recognition per name across formats. In the named individual format, recognition is coded as 0 or 1 for each name. In the paired name format, recognition per name is coded as 0, 0.5, or 1 depending on whether zero, one, or two names in the pair were recognized. Second, we compare the share of Don't know responses among recognized prompts. This tests whether the named individual format induces respondents to provide definite Yes/No answers. Third, we use administrative take-up records for the referenced individuals to compare false positives, \(P(\text{Yes}\mid \text{Truth}=0)\), and false negatives, \(P(\text{No}\mid \text{Truth}=1)\). In the paired name format, the administrative status of the recognized person is identified only when both people in the pair share the same administrative status. We therefore use pairs in which both individuals did not deworm to compute false positives and pairs in which both individuals dewormed to compute false negatives.

\paragraph{Results.}
Table~\ref{tab:demand_checks} summarizes the results. The paired name format increases reported recognition per name by 11.4 percentage points, from 47.8 to 59.3 percent (\(p<0.001\)), consistent with the anonymized format reducing privacy concerns. However, it does not reduce Don't know responses: the Don't know share is 20.4 percent in the named individual format and 19.7 percent in the paired name format. Nor do we find statistically significant evidence that the named individual format inflates reports of peers' participation. The false positive rate is slightly lower in the paired name format than in the named individual format, and the difference is not statistically significant. The paired name format instead increases false negatives among true participants. Overall, these checks provide little evidence that the named individual format mechanically reduces Don't know responses or inflates false positives.

\paragraph{Selection on recognition.}
A remaining concern is that respondents might avoid reporting knowledge of non-participants by saying they do not recognize them in the named individual format. Such selection would matter for the main observability analysis because the Yes/No/Don't know question is only asked conditional on reported recognition. Although we cannot observe latent recognition directly, we can test whether administrative take-up predicts reported recognition differently across formats.

Table~\ref{tab:demand_gradient} reports this exercise. In the named individual format, among prompts with observed administrative truth, recognition is 52.5 percent when the referenced individual truly dewormed and 43.4 percent when the referenced individual did not, a difference of 9.1 percentage points. In the paired name format, we can link truth to recognition only in the subset of pairs where both individuals share the same administrative status. In that subset, recognition per name is 64.4 percent in pairs where both individuals dewormed and 54.8 percent in pairs where both did not, a difference of 9.6 percentage points. Thus, the truth recognition gradient is nearly the same across formats. A difference in differences regression yields an interaction estimate of 0.5 percentage points (SE \(=0.016\), \(p=0.754\)). Moreover, the paired name format raises recognition by roughly the same amount for true non-participants and true participants: paired minus named differences of 11.4 and 11.9 percentage points, respectively. These patterns suggest that format driven differences in recognition reporting are not disproportionately concentrated among true non-participants. This mitigates concerns that selection on recognition in the named individual format is biasing the main treatment effects on observability.


\begin{table}[!htbp]
\centering
\scriptsize
\caption{Elicitation Format Checks for Reported Observability}
\label{tab:demand_checks}
\begin{threeparttable}
\input{tables/demand-checks-table}
\floatfoot{\textit{Notes:}
The sample is restricted to no-SMS endline respondents who completed the social knowledge module and recognized at least one name or pair. The named individual format asks about ten named community members one at a time. The paired name format asks about ten pairs and records how many people in each pair are recognized without requiring respondents to identify which name or names they know. ``Recognition per name'' equals the recognized name indicator in the named individual format and the number recognized divided by two in the paired name format. Don't know shares are computed among recognized prompts. For the paired name format, the Yes/No/Don't know comparison uses prompts in which exactly one person in the pair was recognized, so that the response refers to a single known individual. False positive and false negative rates use administrative take-up records. In the paired name format, truth for the ``person you knew'' question is scored only when both people in the pair share the same administrative status: both did not deworm for false positives and both dewormed for false negatives. Differences and \(p\)-values are from prompt level OLS regressions of each outcome on an indicator for the paired name format, with standard errors clustered by respondent.}
\end{threeparttable}
\end{table}

\begin{table}[!htbp]
\centering
\scriptsize
\caption{Recognition Reporting by Administrative Take-up Status}
\label{tab:demand_gradient}
\begin{threeparttable}
\begin{tabular}{lcccc}
\toprule
 & Truth \(=0\) & Truth \(=1\) & Truth gradient & \(p\)-value \\
\midrule
Named individual format & 0.434 & 0.525 & 0.091 &  \\
Paired name format & 0.548 & 0.644 & 0.096 &  \\
\midrule
Paired minus named & 0.114 & 0.119 & 0.005 & 0.754 \\
\bottomrule
\end{tabular}
\floatfoot{\textit{Notes:}
The sample is restricted to no-SMS endline respondents. In the named individual format, recognition is linked directly to the administrative take-up status of the named individual. In the paired name format, recognition per name is the number recognized divided by two. Because the paired name format does not identify which name was recognized when exactly one name is recognized, the paired name rows restrict to pairs where both individuals have the same administrative take-up status. ``Truth \(=0\)'' denotes named individuals or same status pairs in which the relevant individual or individuals did not deworm. ``Truth \(=1\)'' denotes named individuals or same status pairs in which the relevant individual or individuals dewormed. The truth gradient is the difference in recognition between Truth \(=1\) and Truth \(=0\). The final row reports the difference between paired name and named individual formats. The \(p\)-value is from the interaction term in a prompt level regression of recognition per name on an indicator for the paired name format, the administrative truth indicator, and their interaction, with standard errors clustered by respondent. The interaction estimate is 0.005 (SE \(=0.016\)).}
\end{threeparttable}
\end{table}

\clearpage

\FloatBarrier
\newpage

\subsection{Ruling Out Distance Varying Private Valuation}
\label{sec:mrs}

The reduced form evidence raises a possible private utility alternative to the social image interpretation. If adults at farther distances privately valued the bracelet more relative to the calendar, then the stronger Far effect of the bracelet could reflect distance varying private item values rather than social image. This appendix formalizes the implication of that explanation and relates it to the endline preference evidence.

\paragraph{Latent utility benchmark.}
Let \(B_i,C_i\in\{0,1\}\) indicate assignment to the Bracelet or Calendar arm, with Control given by \(B_i=C_i=0\), and let \(d_i\) be the community-to-PoT distance. Consider a private utility benchmark in which incentives affect take-up only through non-social utility:
\begin{equation}
  y_i^{*}= \alpha
           -\kappa d_i
           +\beta_{\mathrm B} B_i
           +\beta_{\mathrm C} C_i
           +\gamma_{\mathrm B}(B_i d_i)
           +\gamma_{\mathrm C}(C_i d_i)
           +\varepsilon_i,
  \label{eq:latent-private}
\end{equation}
with \(\kappa>0\) and \(y_i=\mathbf 1\{y_i^{*}>0\}\). In this benchmark, the private utility value of the bracelet relative to the calendar at distance \(d\) is
\[
  D(d)
  =
  \frac{\partial y_i^{*}}{\partial B_i}
  -
  \frac{\partial y_i^{*}}{\partial C_i}
  =
  (\beta_{\mathrm B}-\beta_{\mathrm C})
  +
  (\gamma_{\mathrm B}-\gamma_{\mathrm C})d.
\]
Thus, a private valuation explanation for a larger Bracelet effect at farther distances requires the bracelet-calendar private value gap to become more favorable to the bracelet as distance rises:
\[
  \gamma_{\mathrm B}-\gamma_{\mathrm C}>0.
\]
By contrast, if the bracelet-calendar private value gap is distance invariant, then \(\gamma_{\mathrm B}-\gamma_{\mathrm C}=0\), and this channel cannot explain a stronger Bracelet effect in Far communities relative to the Calendar arm.

\paragraph{Empirical implication.}
The preference data provide little support for this implication. In no-item Control communities, where neither item is locally prevalent, the probability that a non-dewormer chooses the bracelet over the calendar is nearly identical in Close and Far communities: 22.1 percent in Close communities and 23.5 percent in Far communities (Table~\ref{tab:pref-cal-gift}). In Bracelet communities, where local prevalence could itself affect private value, bracelet choice among non-dewormers in Far communities is lower than in Far Control communities. These patterns do not indicate that adults at farther distances privately value the bracelet more relative to the calendar.

Consequently, a private utility and distance cost model can rationalize the stronger Far effects of the public signals only by invoking distance varying relative private values that are not observed in the preference data. The calendar placebo, the gift choice evidence, and the relative valuation exercise therefore provide little support for a private value only explanation of the distance by public signal pattern.


\subfile{structural-likelihood}

\clearpage

\subfile{structural-benchmark-recovery}

\clearpage

\subfile{structural-alternative-models}

\clearpage

\begin{table}[ht!]
  \scriptsize
  \caption{Prior Sensitivity of the Structural Social Multiplier at 1.25 km}
  \label{tab:struct-sm-prior-sensitivity-1250m}
  \begin{threeparttable}
  \input{tables/fit105/main-core-multiplier-prior-sensitivity-1250m}
  \floatfoot{\textit{Notes:} Each cell reports the posterior median social multiplier at a distance of 1.25 km, its 95 percent credible interval in square brackets, and the posterior probability of mitigation, $\Pr(M<1)$, on the third line. Posterior probabilities are reported to three decimal places, with values below 0.001 shown as $<0.001$ and values above 0.999 as $>0.999$. Tighter and more diffuse specifications vary the prior for the indicated model component relative to the benchmark; the final two rows vary all listed prior components together.}
  \end{threeparttable}
\end{table}



\clearpage

\subsection{Policymaker PoT Optimization}
\label{sec:oa-appendix}

\subsubsection{Alternative Distance Caps}

\begin{table}[ht!]
  \scriptsize
 \caption{Policymaker PoT Optimization under Alternative Distance Caps}
  \label{tab:optim-post-robust}
  \centering
  \begin{threeparttable}
    \resizebox{\textwidth}{!}{%
      \input{tables/fit105/optim-summ-robust-table}
    }
    \floatfoot{\textit{Notes:}
    Entries report site counts, adult-population-weighted mean take-up, and adult-population-weighted mean one-way distance. Point estimates are posterior medians and parentheses contain 95 percent credible intervals. All caps use the same 1{,}600 posterior draws and, within each draw, the same experimental-Control adult participation target across scenarios. Panel A reports the experimental allocation; site counts and mean distance are fixed, so their intervals are omitted. Panels B--E solve equation~\eqref{eq:oa-problem} using all 1{,}451 candidate sites under the stated distance cap. Control and Bracelet allow social-image returns to vary with distance; fixed-return rows hold the corresponding returns at their 0.5 km values. \textsuperscript{$\dagger$}Without social-image returns (\(\mu=0\)), the target is infeasible in every draw at every cap; these rows report the best feasible allocations. \textsuperscript{$\ddagger$}At 10 km, Control equilibria are undefined in 26 of 1{,}600 draws (1.6 percent); its summaries exclude those draws. All other Control, Bracelet, and fixed-return rows meet the target in every draw.}
  \end{threeparttable}
\end{table}




\FloatBarrier


\subsubsection{Resource-Cost Break-Even Values}

We apply resource-cost accounting to the saved benchmark allocations at the
3.5 km cap. Let $S_z$ denote sites, $Q_z$ expected participants, and $K_z$
expected participant-kilometres of round-trip travel under regime $z$.
At fixed site cost $C_F$ and travel value $C_T$, accounted costs are
\[
C_z = C_F S_z + \mathbf{1}\{z=\mathrm{Bracelet}\}(0.20)Q_z + C_T K_z.
\]
Bracelets cost \$0.20 per expected participant. We compute round-trip travel
using twice the assigned one-way distance, weighted by expected participants
in each community. The fixed site cost at which Bracelet breaks even is
\[
C_F^* = \frac{0.20 Q_{\mathrm{Bracelet}} +
 C_T(K_{\mathrm{Bracelet}}-K_{\mathrm{Control}})}
 {S_{\mathrm{Control}}-S_{\mathrm{Bracelet}}}.
\]
The benchmark saves sites in every retained draw. Costs are evaluated after
minimizing sites subject to adult participation targets; the allocations are
not reoptimized to minimize resource costs.

\begin{table}[htbp]
  \centering
  \small
  \caption{Break-Even Resource Costs for the Bracelet Allocation}
  \label{tab:policy-break-even}
  \begin{threeparttable}
    \resizebox{\textwidth}{!}{\input{tables/fit105/policy-break-even-costs}}
    \floatfoot{\textit{Notes:} Results use the 1{,}600 benchmark posterior
    draws and their saved adult-population-weighted allocations. Break-even
    values are posterior medians, with central 95 percent credible intervals
    in parentheses. Travel values are dollars per participant round-trip
    kilometre. Probabilities give the share of draws in which accounted costs
    are lower under Bracelet than Control at the stated fixed site cost.
    Values are rounded to three decimal places. Both regimes preserve the
    same draw-specific experimental-Control adult participation target.}
  \end{threeparttable}
\end{table}
\FloatBarrier

\subsection{Closed Form for the Type Inference Term \texorpdfstring{\(\Delta(w)\)}{Delta(w)}}
\label{delta-deriv}
\input{appendix-files/analytical-delta.tex}


\end{document}
